\documentclass[a4paper]{article}
\usepackage{parskip}
\usepackage{lipsum}

\def\nterm {Spring}
\def\nyear {2026}
\def\ncourse {Probability}

\newcommand{\convd}{\Rightarrow}
\newcommand{\convw}{\xrightarrow{w}}
\newcommand{\convp}{\xrightarrow{p}}

\input{../header}
\author{Notes taken by \nauthor\vspace{5pt}\\
Lectures by Konstantin Tikhomirov\vspace{5pt}\\
Carnegie Mellon University}

\newcommand{\TODO}{\textcolor{red}{\textbf{*** TO-DO ***}}}

\begin{document}
\maketitle

\tableofcontents

\section{Measure theory review}

\subsection{Measurable space and mapping}

\begin{defi}[$\sigma$-field]
  A collection of subsets $\Sigma \subset 2^{\Omega}$ is a $\sigma$-field if
  \begin{itemize}
    \item $\emptyset \in \Sigma$.
    \item If $A \in \Sigma$, then $A^c \in \Sigma$.
    \item If $\seqinfi{A_i} \subset \Sigma$, then $\cupinfi A_i \in \Sigma$.
  \end{itemize}
  The pair $(\Omega, \Sigma)$ is called a measurable space.
\end{defi}

\begin{defi}[atom]
  Let $\Sigma$ be a $\sigma$-field. Say $A \in \Sigma$ is an atom
  if for all $B \in \Sigma$ either $A \subset B$ or $A \cap B = \emptyset$.
\end{defi}

\begin{prop}
  For all $\omega \in \Omega$, there exists atom $A \in \Sigma$ containing
  $\omega$ if $\Omega$ is finite or countable.
\end{prop}

\begin{proof}
  Define $\tilde{A} = \bigcap \left\{ B \in \Sigma :
    \omega \in B \right\}$. We can check that $\tilde{A} \in \Sigma$
  and $\tilde{A}$ is an atom containing $\omega$.

\end{proof}

\begin{cor}
  If $\Omega$ is finite or countable, there exists a partition
  $\Omega = \bigsqcup_i \Omega_i$,
  where each $\Omega_i$ is an atom of $\Sigma$. With this partition,
  $\Sigma$ is just the power set with respect to $\left\{ \Omega_i \right\}_i$.
\end{cor}

\begin{defi}
  If $F \subset 2^\Omega$, then the $\sigma$-field generated by $F$ is the
  smallest $\sigma$-field containing all elements of $F$.
  Write this $\sigma$-field as $\sigma (F)$.
\end{defi}

\begin{eg}
  Let $\Omega = \left\{ 1, 2, 3, 4, 5 \right\}$ and
  $F = \left\{ \left\{ 2,3 \right\}, \left\{ 3,4 \right\} \right\}$.
  Construct $\sigma$-field $\Sigma$ generated by $F$.
  $\Sigma$ is all possible union of sets from the collection
  $\left\{ \{2\}, \{3\}, \{4\}, \{1, 5\}\right\}$.

\end{eg}

\begin{defi}[measurable mapping]
  Given two measurable spaces $(\Omega, \Sigma)$ and $(\tilde{\Omega},
    \tilde{\Sigma})$. Then $f : \Omega \to \tilde{\Omega}$ is measurable
  if $f^{-1} (B) \in \Sigma$ for all $B \in \tilde{\Sigma}$.
\end{defi}

\begin{defi}[Borel $\sigma$-field]
  Let $(T, \tau)$ be a topological space. Then the Borel $\sigma$-field
  $\B(T, \tau)$ is defined as the smallest $\sigma$-field containing
  all open sets.
\end{defi}

\begin{defi}[product measurable space]
  Given two measurable spaces $(\Omega, \Sigma)$ and
  $(\tilde{\Omega}, \tilde{\Sigma})$. We can define the product
  measurable space as follows: let the ground set be
  $\Omega \times \tilde{\Omega}$, and let
  $\Sigma \otimes \tilde{\Sigma}$ be the smallest $\sigma$-field
  containing all rectangles $B \times \tilde{B}$ where
  $B \in \Sigma$ and $\tilde{B} \in \tilde{\Sigma}$.

  More generally,
  let $\Lambda$ be an index set and $(\Omega_\lambda, \Sigma_\lambda)_{\lambda
      \in \Lambda}$. Define the product $\sigma$-field $\bigotimes_{\lambda
      \in \Lambda} \Sigma_\lambda$ be
  the smallest $\sigma$-field containing all elements
  in the form of $\prod_{\lambda \in \Lambda} B_\lambda$ where
  $B_\lambda \in \Sigma_\lambda$ and $B_\lambda = \Omega_\lambda$ for all
  but countably many indices.
\end{defi}

\begin{prop}
  Let $\left( \Omega_i, \Sigma_i \right)_{i=1}^n$ be measurable
  spaces and $\left( \prod_{i=1}^n \Omega_i, \bigotimes_{i=1}^n \Sigma_i
    \right)$ be the product space. Let $(\Omega, \Sigma)$ be the domain
  and $f = (f_1, \dots, f_n) : (\Omega, \Sigma) \to (\prod_{i=1}^n
    \Omega_i, \bigotimes_{i=1}^n \Sigma_i)$. Suppose $f$ is measurable,
  then every coordinate projection $f_i : \Omega \to \Omega_i$ is
  measurable.

  This is also true for arbitrary index set.
\end{prop}

\begin{prop}
  If $f : (\Omega, \Sigma) \to (\Omega_f, \Sigma_f)$ and
  $g : (\Omega, \Sigma) \to (\Omega_g, \Sigma_g)$, then the concatenation
  $(f, g)$ is measurable w.r.t. the product space
  $(\Omega_f \times \Omega_g, \Sigma_f \otimes \Sigma_g)$.
\end{prop}

\begin{proof}
  Let $A \times B$ be such that $A \in \Sigma_f$ and $B \in \Sigma_g$.
  Then the preimage
  \[
    (f, g)^{-1} (A \times B) = f^{-1} (A) \cap g^{-1} (B) \in \Sigma.
  \]
  By definition, the product $\sigma$-field is generated by rectangles,
  so the proof is complete.

\end{proof}

\subsection{Measure space}
\begin{defi}[measure]
  Let $(\Omega, \Sigma)$ be a measurable space. Then $\mu : \Sigma \to
    [0, \infty]$ is a measure if
  \begin{itemize}
    \item $\mu (\emptyset) = 0$.
    \item If $A_i \in \Sigma$ is pairwise disjoint then
          $\mu \left( \cupinfi A_i \right) = \suminfi \mu (A_i)$.
  \end{itemize}
\end{defi}

\begin{prop}[continuity of measure]
  If $A_1 \subset A_2 \subset \dots$ is a nested sequence of elements
  of $\Sigma$ and $\mu$ be any measure on $(\Omega, \Sigma)$. Then
  \[
    \mu \left( \cupinfi A_i \right) = \lim_{i \to \infty} \mu (A_i).
  \]

  If $A_1 \supset A_2 \supset \dots$ is a nested sequence of elements
  of $\Sigma$ and $\mu (A_n) < \infty$ for some $n$. Then
  \[
    \mu \left( \capinfi A_i \right) = \lim_{i \to \infty} \mu (A_i).
  \]
\end{prop}

\begin{defi}
  Let $(\Omega, \Sigma, \mu)$ be a measure space.

  Say $\mu$ is $\sigma$-finite
  if there is a representation $\Omega = \cupinfi \Omega_i$ where
  $\Omega_i \in \Sigma$ and $\mu (\Omega_i) < \infty$.

  Say $\mu$ is a probability measure if $\mu(\Omega) = 1$.
\end{defi}

\begin{defi}[completion of measure space]
  Let $(\Omega, \Sigma, \mu)$ be a measure space. Let
  \[
    \tilde{\Sigma} = \left\{ A \cup B : A \in \Sigma,
    B \subset \Omega, \text{there exists $C \in \Sigma$ with
  $\mu(C) = 0$ and $B \subset C$} \right\}.
  \]
  We can check $\tilde{\Sigma}$ is a $\sigma$-field. If $\tilde{\mu}$
  is a measure on $(\Omega, \tilde{\Sigma})$ which agrees with $\mu$
  on $\Sigma$, then $(\Omega, \tilde{\Sigma}, \tilde{\mu})$ is called
  a completion of $(\Omega, \Sigma, \mu)$.
\end{defi}

\subsection{$\pi$-$\lambda$ theorem}

\begin{defi}[$\pi$-system]
  Let $\Omega$ be a set and $\P$ be a collection of subsets of $\Omega$.
  Then $\P$ is a $\pi$-system if it is closed with respect to
  taking finite intersections. That is, $A, B \in \P$ implies $A \cap B
    \in \P$.
\end{defi}

\begin{eg}
  On the real line $\R$, both $\P_1 = \left\{ (a, b) : a < b \right\}$
  and $\P_2 = \left\{ (-\infty, a] : a \in \R \right\}$ are
  $\pi$-systems.

\end{eg}

\begin{defi}[$\lambda$-system]
  Let $\Omega$ be a set and $\L$ be a collection of subsets of $\Omega$.
  Say $\L$ is a $\lambda$-system if
  \begin{itemize}
    \item $\emptyset \in \L$.
    \item $A \in \L$ implies $A^c \in \L$.
    \item for all countable collection of disjoint elements $A_i \in \L$,
          we have $\cupinfi A_i \in \L$.
  \end{itemize}

  For an alternative definition, say $\L$ is a $\lambda$-system if
  \begin{itemize}
    \item $\Omega \in \L$.
    \item If $A, B \in \L$ and $A \subset B$, then $B \setminus A \in \L$.
    \item If $A_n \in \L$ and $A_n \uparrow A$, then $A \in \L$.
  \end{itemize}
\end{defi}

\begin{thm}[$\pi$-$\lambda$ theorem]
  Let $\Omega$ be a set, $\P$ be a $\pi$-system and $\L$ be a $\lambda$-system.
  Also suppose $\P \subset \L$, then $\sigma (\P) \subset \L$.
\end{thm}

\begin{proof}
  Let $\ell (\P)$ be the smallest $\lambda$-system on $\Omega$ containing
  $\P$. The goal is to show that $\ell (\P)$ is a $\sigma$-field.
  We need to show that if $A_i \in \ell (\P)$ for $1 \leq i < \infty$,
  then $\cupinfi A_i \in \ell (\P)$. Note that
  \[
    \cupinfi A_i = \cupinfi \left( A_i \setminus \cupj^{i - 1} A_i \right),
  \]
  so it suffices to show that $A, B \in \ell (\P)$ implies $A \cap B \in
    \ell (\P)$.

  For $A \in \ell(\P)$ we define
  \[
    W_A = \left\{ B \subset \Omega : A \cap B \in \ell (\P) \right\}.
  \]
  It can be directly verified that $W_A$ is a $\lambda$-system.

  Take $A \in \P$, then for any $B \in \P$ we have $A \cap B \in \P
    \subset \ell (\P)$. Hence, $\P \subset W_A$ and thus $\ell (\P) \subset W_A$
  for all $A \in \P$, as $\ell (\P)$ is the smallest $\lambda$-system on
  $\Omega$ containing $\P$.
  Now take $A \in \ell (\P)$, we have $A \in W_B$
  for all $B \in \P$. It follows that $A \cap B \in \ell (\P)$ and thus
  $B \in W_A$. Hence similarly
  $\ell (\P) \subset W_A$ for all $A \in \ell (\P)$.

  Now for any pair $B, C \in \ell (\P)$, we have $C \in W_B$ and thus
  $B \cap C \in \ell (\P)$. This completes the proof.

\end{proof}

\subsection{Extension theorems}

\begin{defi}[semi-field]
  A collection of subsets $S \subset 2^{\Omega}$ is a semi-field if
  \begin{itemize}
    \item $\emptyset \in S$ and $\Omega \in S$.
    \item $A, B \in S$ implies $A \cap B \in S$.
    \item If $A \in S$, then $A^c$ is a finite disjoint union of sets in $S$.
  \end{itemize}
\end{defi}

\begin{thm}[Caratheorody's extension theorem]
  Let $S$ be a semi-field and let $\mu$ be a non-negative function on $S$
  satisfying:
  \begin{itemize}
    \item $\mu (\emptyset) = 0$.
    \item If $A_1, \dots, A_n$ are disjoint and $\cupi^n A_i \in S$, then
          $\mu (\cupi^n A_i) = \sumi^n \mu (A_i)$.
    \item If $A_1, A_2, \dots$ are such that $\cupinfi A_i \in S$, then
          $\mu (\cupinfi A_i) \leq \suminfi \mu (A_i)$.
  \end{itemize}
  Then $\mu$ admits a unique extension $\bar{\mu}$ which is a
  measure on $\bar{S}$, the field (algebra) generated by $S$. Moreover,
  if $\bar{\mu}$ is $\sigma$-finite then $\bar{\mu}$ admits a unique
  extension $\tilde{\mu}$ to $\sigma (S)$.
\end{thm}

\begin{notation}
  Let $T$ be any set. Write
  \[
    \R^T = \left\{ \left( \omega_t \right)_{t \in T} :
    \omega_t \in \R \right\}.
  \]
  Also write $\rcal^T$ as the $\sigma$-field generated by rectangles
  of the form $\prod_{t \in T} I_t$, where for each $t \in T$,
  $I_t$ is either a semi-open interval of the form $(a, b]$
  with $a < b$ or $I_t = \R$, and $I_t = \R$ for all but finitely
  many $t \in T$.
\end{notation}

\begin{thm}[Kolmogorov's extension theorem]
  For each finite non-empty subset $J \subset T$,
  let $\mu_J$ be a Borel probability measure in $\R^J$,
  and assume that the measures $\left( \mu_J \right)_{J \subset T,
      \abs{J} < \infty}$ are
  compatible, in the sense that whenever $J_1 \subset J_2 \subset T$
  with $0 \leq \abs{J_1} \leq \abs{J_2} < \infty$,
  $I_j \subset \R$ with $j \in J_1$ are Borel subsets of $\R$, and
  \[
    \tilde{I_j} = \begin{cases}
      I_j & (j \in J_1)                \\
      \R  & (j \in J_2 \setminus J_1),
    \end{cases}
  \]
  one has
  \[
    \mu_{J_2} \left( \prod_{j \in J_2} \tilde{I_{j}} \right)
    = \mu_{J_1} \left( \prod_{j \in J_1} I_j \right).
  \]
  Then there exists a unique probability measure $\mu$ on
  $(\R^T, \rcal^T)$ consistent with $\left( \mu_J \right)_{J \subset T,
      \abs{J} < \infty}$. That is, one has
  \[
    \mu \left( \prod_{t \in T} I_t \right)
    = \mu_J \left( \prod_{j \in J} I_j \right)
  \]
  whenever $J \subset T$ with $\abs{J} < \infty$ and $I_t = \R$
  for all $t \notin J$.
\end{thm}

\subsection{Lebesgue Integration}

Here we provide a proof for dominated convergence theorem that uses
the truncation technique, which will be a useful technique later
in the course.

\begin{thm}[dominated convergence theorem]
  Let $\seqinfn{f_n}$ be a sequence of measurable
  functions on $(\Omega, \Sigma, \mu)$ and $g \geq 0$ be another measurable
  function. Suppose
  \begin{enumerate}
    \item $\int g \d \mu < \infty$.
    \item $\abs{f_n} (\omega) \leq g (\omega)$ for all $\omega \in \Omega$
          and $n \geq 1$.
    \item $f_n \to f$ pointwise.
  \end{enumerate}
  Then
  \[
    \lim_{n \to \infty} \int f_n \d \mu = \int f \d \mu.
  \]
\end{thm}

\begin{proof}
  \textbf{Claim 1.}
  If $h$ is a function on $(\Omega, \Sigma, \mu)$ with
  $h \geq 0$ and $\int h \d \mu < \infty$. Let $\seqinfn{A_n}$
  be any sequence of elements of $\Sigma$ with $\mu (A_n) \to 0$.
  Then
  \[
    \int_{A_n} h \d \mu \to 0.
  \]
  Proof of claim. WLOG assume $\mu(A_n) \leq 2^{-n}$ for all $n$.
  Define $h_n = h \ind_{\bigcup_{i=n}^\infty A_i}$. We then have
  \begin{enumerate}
    \item The sequence $\seqinfn{h_n}$ is monotone.
    \item $h_n$ converges to $0$ almost everywhere.
  \end{enumerate}
  Monotone convergence theorem then implies $\lim_{n \to \infty}
    \int h_n \d \mu = 0$. Meanwhile,
  \[
    0 \leq \int_{A_n} h \d \mu \leq \int h_n \d \mu,
  \]
  and the proof is complete.

  \textbf{Claim 2.}
  Suppose $h \geq 0$ and $\int h \d \mu < \infty$.
  Let $\seqinfn{\epsilon_n}$ be a sequence of strictly positive numbers
  converging to zero. Define
  \[
    B_n = \left\{ \omega \in \Omega : h (\omega) \leq \epsilon_n \right\}
    \in \Sigma.
  \]
  Then
  \[
    \int_{B_n} h \d \mu \to 0.
  \]
  Proof of this claim is left as an exercise.

  Now we prove the theorem. Fix $\epsilon > 0$. By the previous two
  claims, there exists $M > 0$ and $\delta > 0$ such that
  \[
    \int_{\left\{ g \geq M \right\}} g \d \mu < \epsilon,
    \quad
    \int_{\left\{ g \leq \delta \right\}} g \d \mu < \epsilon.
  \]
  Let $U = \left\{ \omega : \delta < g(\omega) < M \right\}$.
  Since $g$ is integrable, $\mu (U) < \infty$.
  For $\omega \in U$, let $n_\epsilon (\omega)$ be the smallest
  index such that $n \geq n_\epsilon (\omega)$ implies
  $\abs{f_n (\omega) - f(\omega)} \leq \epsilon \mu(U)^{-1}$.
  It follows that there exists $N$ such that
  \[
    \mu \left( \left\{ \omega \in U: n_\epsilon (\omega) > N \right\}  \right)
    \leq \frac{\epsilon}{M}.
  \]
  Then, for $n \geq N$, we have
  \begin{align*}
    \abs{\int_U (f_n - f) \d \mu}
    \leq \int_{n_\epsilon (\omega) \leq N} \abs{f_n - f} \d \mu
    + \int_{n_\epsilon (\omega) > N} \abs{f_n - f} \d \mu
    \leq 3 \epsilon.
  \end{align*}
  Now for $n \geq N$, we have
  \[
    \abs{\int (f_n - f) \d \mu}
    \leq 3 \epsilon + \int_{U^c} \abs{f - f_n} \d \mu
    \leq 3 \epsilon + 2 \int_{U^c} g \d \mu \leq 7 \epsilon.
  \]

\end{proof}

\begin{thm}[Markov-Chebyshev inequality]
  Suppose we have probability measure space $(\Omega, \Sigma, P)$
  and $f \geq 0$. Suppose also $\int f \d P < \infty$. Then
  \[
    P \left( \left\{ \omega : f(\omega) > t \right\} \right)
    \leq \frac{1}{t} \int f \d P.
  \]
  for all $t > 0$.
\end{thm}

\begin{remark}
  Let $1 \leq p < \infty$. Suppose $f : (\Omega, \Sigma, P) \to
    [0, \infty]$ and $\int f^p \d P < \infty$. Then
  \[
    P \left( \left\{ \omega : f (\omega) > t \right\} \right)
    \leq \frac{1}{t^p} \int f^p \d P.
  \]
  for all $t > 0$.
\end{remark}

\begin{remark}
  Suppose $\int e^{\lambda f} \d P < \infty$ for all
  $\lambda \in \R$ and $f : (\Omega, \Sigma, P) \to [0, \infty]$.
  Then
  \[
    P \left( \left\{ \omega : f (\omega) > t \right\} \right)
    \leq \frac{1}{e^{\lambda t}} \int e^{\lambda f} \d P
  \]
  for all $t > 0$ and $\lambda > 0$.
\end{remark}

\begin{thm}[H\"older inequality]
  Let $p, q \in [1, \infty]$ and $p^{-1} + q^{-1} = 1$.
  Let $(\Omega, \Sigma, \mu)$ be a probability space. For any
  measurable functions $f, g$, we have
  \[
    \int \abs{f g} \d P \leq \left( \int \abs{f}^p \d P \right)^{1 / p}
    \left( \int \abs{g}^q \d P \right)^{1 / q}.
  \]
\end{thm}

\begin{thm}[Jensen's inequality]
  Let $(\Omega, \Sigma, P)$ be a probability space and $f$ be
  integrable. Let $\phi : \bar{\R} \to \bar{\R}$ be convex
  and suppose $\phi (\infty) = \lim_{x \to \infty} \phi(x)$
  and $\phi (-\infty) = \lim_{x \to -\infty} \phi(x)$. Then
  \[
    \phi \left( \int f \d P \right) \leq \int \phi (f) \d P.
  \]
\end{thm}

\subsection{Product measures and Fubini theorem}
Let $(\Omega_1, \Sigma_1, \mu_1)$, $(\Omega_2, \Sigma_2, \mu_2)$
be $\sigma$-finite measure spaces. We already defined the product
$\Sigma_1 \otimes \Sigma_2$. To define a product measure, we first
consider the algebra of rectangles
\[
  S = \left\{ A \in \Sigma_1 \otimes \Sigma_2 : A = A_1 \times A_2
  \text{ for some $A_1 \in \Sigma_1, A_2 \in \Sigma_2$} \right\}.
\]
Then we can define $\mu = \mu_1 \times \mu_2$ on $S$ by
\[
  \mu (A) = \mu_1 (A_1) \mu_2 (A_2)
\]
for $A = A_1 \times A_2$. We can check that the definition is
self-consistent. That is, if $A = A_1 \times A_2$ is a countable
union of disjoint rectangles
$\seqinfj{A_1^{(j)} \times A_2^{(j)}}$, we have
\[
  \mu (A_1 \times A_2) = \suminfj \mu ( A_1^{(j)} \times A_2^{(j)} ).
\]
This can be verified with monotone convergence theorem.
Now $\mu$ is a premeasure and can be uniquely
extended to $\Sigma_1 \otimes \Sigma_2$.

\begin{thm}[Fubini-Tonelli]
  Let $(\Omega_1, \Sigma_1, \mu_1)$, $(\Omega_2, \Sigma_2, \mu_2)$
  be $\sigma$-finite measure spaces and let $(\Omega, \Sigma, \mu)$
  be the product space. Suppose $f$ is measurable on the product space.
  Suppose either $f$ is non-negative or $\int_\Omega \abs{f} \d \mu <
    \infty$. Then
  \begin{itemize}
    \item $y \mapsto f(x, y)$ is $\Sigma_2$ measurable for all $x \in \Omega_1$.
    \item $x \mapsto \int_{\Omega_2} f(x, y) \d \mu_2 (y)$ is
          $\Sigma_1$ measurable.
    \item We have
          \[
            \int_{\Omega_1} \int_{\Omega_2} f(x, y) \d \mu_2 (y) \d \mu_1 (x)
            = \int_\Omega f(x, y) \d \mu(x, y).
          \]
  \end{itemize}
\end{thm}

\begin{proof}
  First suppose $f = \ind_{A}$ for $A \in \Sigma$.
  Also suppose $\mu_1, \mu_2$ are finite. Define section
  \[
    A_x = \left\{ y \in \Omega_2 : (x, y) \in A \right\}.
  \]
  The goal is to show that $A_x \in \Sigma_2$ for all $x \in \Omega_1$.
  Define a family of sets
  \[
    \fcal_x = \left\{ B \in \Sigma : \text{$B_x$ is $\Sigma_2$-measurable} \right\}.
  \]
  It can be verified that $\fcal_x$ is a $\sigma$-field for all $x \in \Omega_1$.
  Also, $\fcal_x$ contains all rectangles and thus $\Sigma \subset \fcal_x$.
  Hence, we have shown that $y \mapsto \ind_A (x, y) = \ind_{A_x} (y)$
  is measurable for all $x \in \Omega_1$.

  Next we show $x \mapsto \mu_2 (A_x)$ is measurable
  and its integral over $\Omega_1$ is equal to $\mu (A)$.
  Define
  \[
    \ucal = \left\{ B \in \Sigma : \text{$x \mapsto \mu_2(B_x)$ is
  $\Sigma_1$-measurable and $\int_{\Omega_1} \mu_2(B_x) \d \mu_1 = \mu (B)$} \right\}
  \]
  It can be verified that $\ucal$ is a $\lambda$-system.
  Note that $\ucal$ also contains all rectangles in $\Sigma$.
  It follows that $\ucal = \Sigma$ and the proof for
  indicator functions are complete.

  Then use linearity to extend to simple functions, and use monotone
  convergence theorem to prove the statement for non-negative functions.
  For the case where $f$ is integrable, consider the positive and
  negative part about $f$ to complete the proof.

\end{proof}

\section{Probability theory basics}

\subsection{Distributions and densities}

\begin{defi}
  Let $F : \R \to [0, 1]$. Suppose $F$ is
  \begin{itemize}
    \item right-continuous.
    \item non-decreasing.
    \item $\lim_{t \to -\infty} F(t) = 0$ and
          $\lim_{t \to \infty} f(t) = 1$.
  \end{itemize}
  Then $F$ is a cumulative distribution function (CDF).
\end{defi}

\begin{remark}
  If we want to define CDF in $\R^2$ then the axioms are
  \begin{itemize}
    \item right-continuous: $F(\tilde{s}, \tilde{t})
            \to F(s, t)$ as $t \downarrow \tilde{t}$ and
          $s \downarrow \tilde{s}$.
    \item coordinate-wise non-decreasing.
    \item $\lim_{s, t \to \infty} F(s, t) = 1$,
          $\lim_{s \to -\infty} F(s, t) = 0$ for any $t$, and
          $\lim_{t \to - \infty} F (s, t) = 0$ for any $s$.
    \item For a rectangle with bottom left vertex $(a_1, a_2)$
          and top right vertex $(b_1, b_2)$,
          \[
            F(b_1, b_2) - F(b_1, a_2) - F(a_1, b_2) + F(a_1, a_2) \geq 0.
          \]
  \end{itemize}
\end{remark}

Now we can connect the notion of CDF with randomness.

Suppose $X$-random real-valued variable on $(\Omega, \Sigma, P)$
that is almost everywhere finite. Define
\[
  F_X (t) = P \left( X(\omega) \leq t \right)
\]
for $- \infty < t < \infty$. It can be verified that
$F_X$ is a CDF.

Conversely, for any CDF $F$, there exists a probability space
$(\Omega, \Sigma, P)$ and a real valued random variable
on $(\Omega, \Sigma, P)$ with CDF $F$.

\begin{defi}
  If $X$ is random variable on $(\Omega, \Sigma, P)$ real valued and a.e.\ finite.
  Then we can define the induced Borel probability measure
  $\mu_X$ on $(\R, \bcal_\R)$ by
  \[
    \mu_X (B) := P \left( X \in B \right)
  \]
  for all $B \in \bcal_\R$.
\end{defi}

Now suppose $\mu$ is any Borel probability measure on $\R$. Consider
probability space $(\R, \bcal_\R, \mu)$ and formal identity mapping
$\id$ on $\R$. Then $\mu_{\id} \equiv \mu$.

\begin{thm}
  There is a one-to-one correspondence between the family of CDFs
  and the family of Borel probability measure on $\R$.
\end{thm}

\begin{proof}
  For any Borel probability measure $\mu$,
  $F_\mu (t) = \mu((-\infty, t])$
  is a valid CDF.

  Conversely, for any CDF $F$, there exists unique probability
  measure $\mu_F$ on $\R$ such that $\mu_F ((-\infty, t]) = F(t)$
  for all $-\infty < t < \infty$. This is a corollary of Caratheorody
  extension theorem. For detailed proof see notes or textbook.

\end{proof}

\begin{remark}
  Suppose $X = (X_1, X_2)$ is a random vector in $\R^2$. We can
  define
  \[
    F_X (s, t) = P \left( X_1 \leq s, X_2 \leq t \right).
  \]
  Corresponding results are also true.
\end{remark}

\begin{defi}[Probability mass function]
  Let $(\Omega, \Sigma, P)$ be a probability space and $X: \Omega
    \to \R$ be random variable. Suppose there exists $S \subset \R$ countable
  such that $P \left( X \in S \right) = 1$. We can define
  the probability mass function (PMF) $f_X$ via
  \[
    f_X (t) = P \left( X = t \right)
  \]
  for $t \in \R$. Due to the restriction, this gives complete description
  of the distribution, and we can construct CDF $F_X$ via
  \[
    F_X (t) = \sum_{s \leq t} f_X (s).
  \]
  This sum makes sense since the $f_X (s) = 0$ for all but countably many
  $s$.
  Conversely, we can also reconstruct $f_X$ from a CDF $F_X$.
\end{defi}

\begin{defi}[Probability density function]
  Suppose $F$ is a CDF which is absolutely continuous. That is,
  there exists Borel measurable non-negative function $\rho$ on $\R$
  such that
  \[
    F(t) = \int_{-\infty}^t \rho (s) \d s
  \]
  for all $-\infty < t < \infty$.
  This implies $F$ is almost everywhere differentiable and the derivative
  is $\rho$. In this case, say $\rho$ is the density function.

  If random variable $X$ is such that $F_X$ is absolutely continuous, then the
  corresponding $\rho_X$ is the probability density function for $X$.
\end{defi}

\begin{remark}
  Recall that a Borel $\sigma$-finite measure $\mu$ on the real line
  is absolutely continuous w.r.t the Lebesgue measure $m$ on $\R$ if
  $\mu (A) = 0$ whenever $A \in \bcal_\R$ is Lebesgue null.
  In this case, Randon-Nikodym theorem implies existence of non-negative
  Borel measurable function $f$ such that $\mu (A) = \int_A f \d m$.
\end{remark}

\begin{thm}
  Suppose $X$ RV on $(\Omega, \Sigma, P)$ is real-valued and a.e.\ finite.
  The following are equivalent:
  \begin{enumerate}
    \item $F_X$ is absolutely continuous.
    \item $\mu_X$ is absolutely continuous w.r.t. Lebesgue measure.
  \end{enumerate}
  Moreover, $\rho_X$ is also the derivative of $\mu_X$ w.r.t. Lebesgue
  measure. That is, for any $A \in \bcal_\R$,
  \[
    \mu_X (A) = \int_A \rho_X (t) \d t.
  \]
\end{thm}

\subsection{Independence}

\begin{defi}
  Say two events $A, B \in \Sigma$ are independent if
  $P (A \cap B) = P (A) P (B)$.
\end{defi}

It is easy to verify that $A, B$ are independent implies
$A^c, B$ are independent.

\begin{remark}
  Suppose $P (B) > 0$, then the conditional probability of $A$
  given $B$ is defined as
  \[
    P (A | B) = \frac{P (A \cap B)}{P (B)}.
  \]
  Then, independence of $A$ and $B$ is equivalent to
  $P (A | B) = P (A)$.
\end{remark}

\begin{defi}
  Let $A_1, \dots, A_n$ be events. Say they are mutually independent
  if for any $\emptyset \neq I \subset [n]$, we have
  \[
    P \left( \bigcap_{i \in I} A_i \right) = \prod_{i \in I} P (A_i).
  \]
  This is equivalent to saying that for every $2 \leq i \leq n$,
  the event $A_i$ is independent from any event generated by
  $A_1, \dots, A_{i - 1}$, or $A_i$ is independent from
  $\sigma \left( A_1, \dots, A_{i - 1} \right)$.
\end{defi}

\begin{remark}
  The events $A_1, \dots, A_n$ are called $k$-wise independent if
  any $k$-subset of the events are mutually independent. For $k < n$,
  this notion is strictly weaker than mutual independence of all
  $n$ events. As an example, consider $P$ to be the uniform distribution
  on $\left\{ 1, \dots, 4 \right\}$. Let $A_1 = \left\{ 1, 2 \right\}$,
  $A_2 = \left\{ 1, 3 \right\}$, and $A_3 = \left\{ 2, 3 \right\}$.
  Then they are pairwise independent but not mutually independent.
\end{remark}

\begin{defi}
  A collection of events $\{A_\lambda\}_{\lambda \in \Lambda}$
  on $(\Omega, \Sigma, P)$ are mutually independent if any finite subset
  of events are mutually independent.
\end{defi}

\begin{defi}
  Let $(\Omega, \Sigma, P)$ be a probability space. Two
  $\sigma$-subfields are independent if for any $A \in \Sigma_1$
  and $B \in \Sigma_2$, $A, B$ are independent.
\end{defi}

\begin{defi}
  Let $(\Omega, \Sigma, P)$ be a probability space and
  $X, Y$ be two real-valued random variables. Say $X$ and $Y$
  are independent if
  \[
    P \left( X \in A, Y \in B \right) = P \left( X \in A \right)
    P \left( Y \in B \right)
  \]
  for any $A, B \in \bcal_\R$.

  Equivalently, let $\Sigma_X, \Sigma_Y$
  be the $\sigma$-field generated by $X$ and $Y$. Then independence of
  $X$ and $Y$ is equivalent to independence of $\Sigma_X$ and $\Sigma_Y$.
\end{defi}

Now we explore how this connect with product structure.

\begin{prop}
  Let $(\Omega_1, \Sigma_1, \Pr_1)$ and $(\Omega_2, \Sigma_2, \Pr_2)$
  be two probability spaces and let $(\Omega, \Sigma, P)$ be the
  product space. Let $X$ and $Y$ be two random variables on
  $(\Omega, \Sigma, P)$. Suppose there exists some measurable
  functions such that $X (\omega_1, \omega_2) = g(\omega_1)$,
  and $Y(\omega_1, \omega_2) = h (\omega_2)$. Then $X$ and $Y$ are
  independent.
\end{prop}

\begin{proof}
  Let $A, B \in \bcal_\R$. Then
  \begin{align*}
    P \left( X \in A, Y \in B \right)
     & = P \left( (\omega_1, \omega_2) : \omega_1 \in g^{-1} (A),
    \omega_2 \in h^{-1} (B) \right)                               \\
     & = P \left( \left\{ \omega_1 \in g^{-1} (A) \right\} \times
    \left\{ \omega_2 \in h^{-1} (B) \right\} \right)              \\
     & = \Pr_1 \left( \omega_1 \in g^{-1} (A) \right)
    \Pr_2 \left( \omega_2 \in h^{-1} (B) \right).
  \end{align*}
  However,
  \begin{align*}
    \Pr_1 \left( \omega_1 \in g^{-1} (A) \right)
     & = P \left( (\omega_1, \omega_2) : \omega_1 \in g^{-1} (A),
    \omega_2 \in \Omega_2 \right)                                 \\
     & = P \left( X \in A \right),
  \end{align*}
  and similarly for $Y$.

\end{proof}

\begin{remark}
  Let $(\Omega, \Sigma, P)$ be a probability space and suppose
  $X, Y$ be two random variables that are independent and
  a.e.\ finite.
  They then generate two Borel probability measure $\mu_X$ and
  $\mu_Y$ on $\R$.
  Define a product probability space as
  of $(\R^2, \bcal_{\R^2}, \mu_X \times \mu_Y)$. Define
  $\tilde{X} (x, y) = x$ and $\tilde{Y} (x, y) = y$ as random
  variables on the product space. By definition,
  $\tilde{X}$ is equidistributed with $X$. That is, $\mu_{\tilde{X}} = \mu_X$
  and $F_{\tilde{X}} = F_X$. Similarly $\mu_{\tilde{Y}} = \mu_Y$.
  Also, $\tilde{X}, \tilde{Y}$ are independent. Now
  $(X, Y)$ and $(\tilde{X}, \tilde{Y})$ have the same distribution.
\end{remark}

\begin{remark}
  If $X$ and $Y$ are independent, then their joint distribution
  $F_{(X, Y)}$ is uniquely determined by the individual distributions
  of $F_X, F_Y$. Indeed,
  \[
    F_{(X, Y)} (s, t) = F_X (s) F_Y (t).
  \]
\end{remark}

\begin{remark}
  Let $(\Omega, \Sigma, P)$ be a probability space and suppose
  $X, Y$ be two random variables that are independent. Suppose they have
  densities $\rho_X, \rho_Y$, then the distribution density of
  vector $(X, Y)$ is $\rho_{(X, Y)} (s, t) = \rho_X (s) \rho_Y (t)$.
\end{remark}

\begin{remark}
  If $X$ and $Y$ are independent random variable, and $f, g$ are measurable functions.
  Then $f(X)$ and $g(Y)$ are independent as well.
\end{remark}

\begin{remark}
  Given probability space $(\Omega, \Sigma, P)$ and random variable $X$.
  It may not exists another random variable $Y$ that is independent from $X$
  on the same probability space.
  See the following example.
\end{remark}

\begin{eg}
  As an example, consider $([0, 1], \bcal_{[0, 1]}, m)$ and
  $X(\omega) = \omega$, so $X$ is uniform on $[0, 1]$. The goal is
  to construct variable $Y$ such that $Y \sim \Bernoulli (\frac{1}{2})$.

\end{eg}

\begin{proof}
  Let $A \in \bcal_{[0, 1]}$ and $t \in [0, 1]$. Define the density
  of set $A$ at point $t$ to be
  \[
    \lim_{\epsilon \to 0} \frac{m (A \cap [t - \epsilon, t + \epsilon])}
    {2 \epsilon}.
  \]
  The Lebesgue density theorem says this is well defined and
  takes values in $\left\{ 0, 1 \right\}$ for $m$-a.e.\ $t \in [0, 1]$.

  Suppose such $Y$ exists, we can vies $Y$ as an indicator of a set
  $A \subset [0, 1]$ of probability $\frac{1}{2}$. That is,
  $Y = \ind_A$ and $P (A) = \frac{1}{2}$. Choose any point $t \in (0, 1)$
  such that density of $A$ at $t$ is well-defined. WLOG assume
  the density is $1$. Pick $\epsilon > 0$ such that
  $m (A \cap [t - \epsilon, t + \epsilon]) \geq \frac{3}{2} \epsilon$.
  Now,
  \begin{align*}
    m (A \cap [t - \epsilon, t + \epsilon])
     & = P \left( Y = 1, X \in [t - \epsilon, t + \epsilon] \right)         \\
     & = P \left( Y = 1 \right) P \left( X \in [t - \epsilon, t + \epsilon]
    \right)                                                                 \\
     & = \epsilon,
  \end{align*}
  a contradiction.

  Alternatively, we can derive a contradiction using
  $m(A) = P \left( Y = 1, X \in A \right)$.

\end{proof}

\begin{remark}
  The goal is to have statements ``independent'' from the underlying
  probability space.
\end{remark}

\subsection{Convolution}

\begin{defi}[convolution]
  Let $\mu$, $\nu$ be Borel probability measures on $\R$.
  The convolution of $\mu$ and $\nu$ is a probability measure
  on $\R$ such that
  \[
    (\mu * \nu) (S) = \int_{\R^2} \ind_S (x + y) d (\mu \times \nu)
  \]
  for all $S \in \bcal_{\R}$.
\end{defi}

\begin{remark}
  Suppose both $\mu$ and $\nu$ have densities $\rho_\mu$ and
  $\rho_\nu$. That is,
  $\mu \ll m$ and $\nu \ll m$. It follows that
  \begin{align*}
    (\mu * \nu) (S)
     & = \int_\R \left( \int_\R \ind_S (x + y) \rho_\mu (x) \d x \right)
    \rho_\nu (y) \d y                                                          \\
     & = \int_\R \left( \int_S \rho_\mu (w - y) \d w \right) \rho_\nu (y) \d y \\
     & = \int_S \left( \int_\R \rho_\mu (w - y) \rho_\nu (y) \d y \right)
    \d w
  \end{align*}
  Note that this implies $\mu * \nu \ll m$ and
  \[
    \rho_{\mu * \nu} (w) = \int_\R \rho_\mu (w - y) \rho_\nu (y) \d y,
  \]
  which is the convolution of funtion $\rho_\mu$ and $\rho_\nu$.
\end{remark}

\begin{remark}
  If $X$ and $Y$ are two independent random variables and
  $\mu_X$ and $\mu_Y$ are the corresponding induced Borel measure $\R$,
  then $\mu_{X + Y} = \mu_X * \mu_Y$.
\end{remark}

\begin{remark}
  Suppose $X$ and $Y$ are independent and their CDF is $F_X$ and $F_Y$,
  then
  \[
    F_{X + Y} (t) = \int_\R F_X (t - w) \d \mu_Y (w)
    = \int_\R F_X (t - w) \d F_Y (w).
  \]
\end{remark}

\subsection{Moments}

In this section we explore the computation and basic properties of
moments.

\begin{defi}[moments]
  Let $X$ be a random variable on $(\Omega, \Sigma, P)$.
  The $p$-th absolute moment of $X$ is
  \[
    E \abs{X}^p = \int_\Omega \abs{X}^p \d P.
  \]
  This is always well-defined but can be infinite.

  If $p$ is positive integer, the $p$-th moment of $X$ is
  \[
    E X^p = \int_\Omega X^p \d P
  \]
  whenever it is defined.

  In particular, $E X$ is the mean or expectation, the variance is
  $\var (X) = E (X - E X)^2$ whenever the expectation is defined,
  and the standard deviation is $\sqrt{\var (X)}$.
\end{defi}

\begin{prop}
  Let $X$ be random variable $(\Omega, \Sigma, P)$ and
  $0 < p < q < \infty$. Then,
  \[
    \left( E \abs{X}^p \right)^{1 / p} \leq \left( E \abs{X}^q \right)^{1 / q}.
  \]
\end{prop}

\begin{proof}
  Define $Y = \abs{X}^p$, then we want to show that
  $E Y \leq (E Y^{q / p})^{p / q}$. Note that $t \mapsto \abs{t}^{q / p}$
  is convex. Therefore, by Jensen's inequality,
  \[
    \left( E Y \right)^{q / p}  \leq E \left( Y^{q / p} \right).
  \]

\end{proof}

Now we want to show that moments only depends on the distribution of the
random variable, and it does not carry unnecessary information about the
underlying probability space. To do this, we first show the following
proposition.

\begin{prop}
  Let $g$ be measurable and $X$ a random variable. Then,
  \[
    E \abs{g(X)} = \int_\R \abs{g(t)} \d \mu_X (t).
  \]
  Moreover, if $E \abs{g(X)} < \infty$, then
  \[
    E g(X) = \int_\R g(t) \d \mu_X (t).
  \]
\end{prop}

\begin{cor}
  If $E X$ is well-defined, then
  \[
    E X = \int_\R t \d \mu_X (t).
  \]
  Moreover, if $X$ has distribution density $\rho_X$, then
  $E X = \int_\R t \rho_X (t) \d t$.
\end{cor}

\begin{prop}
  If $X \geq 0$, then
  \[
    E X = \int_0^\infty P \left( X \geq t \right) \d t.
  \]
  In particular, if $X$ is non-negative and integer valued, then
  $E X = \suminfi P \left( X \geq i \right)$.
\end{prop}

\begin{proof}
  With Fubini-Tonelli, we have
  \[
    E X = \int_0^\infty s \d \mu_X (s)
    = \int_0^\infty \int_0^s 1 \d t \d \mu_X (s)
    = \int_0^\infty \int_t^\infty 1 \d \mu_X \d t
    = \int_0^\infty P \left( X \geq t \right) \d t.
  \]

\end{proof}

\begin{prop}
  Let $X_1, \dots, X_n$ are random variables on $(\Omega, \Sigma, P)$.
  Suppose they are either all non-negative or all integrable.
  Suppose also that they are mutually independent. Then,
  \[
    E \left( \prod_{i=1}^n X_i \right) = \prod_{i=1}^n E X_i.
  \]
\end{prop}

\begin{proof}
  It suffices to show the statement for $n = 2$. Define
  independent variables $\tilde{X_1}$ and $\tilde{X_2}$
  on the product space $(\R^2, \bcal_{\R^2}, \mu_{X_1} \times \mu_{X_2})$
  by coordinate projection. We know $\tilde{X_i}$ is equidistributed
  with $X_i$, so
  \[
    E [X_1 X_2] = E [\tilde{X_1} \tilde{X_2}] = E [\tilde{X_1} ] E
      [ \tilde{X_2} ] = E [X_1] E [X_2] ,
  \]
  where in the second equality we used Fubini-Tonelli.

\end{proof}

\begin{remark}
  Recall that whenever $E X$ is finite, $\var X = E (X - E X)^2
    = E X^2 - (E X)^2$. If $X_1, \dots, X_n$ are pairwise independent
  and have well-defined variances, then
  \[
    \var (X_1 + \dots + X_n) = \sumi^n \var X_i.
  \]
\end{remark}

\begin{defi}[moment generating function]
  Let $X$ be a random variable, and $\lambda \in \R$. Define the
  moment generating function of $X$ via
  \[
    M_X (\lambda) = E \exp (\lambda X).
  \]
\end{defi}

Suppose $M_X (\lambda)$ is finite in some neighborhood of $0$. Then,
\[
  M_X (\lambda)
  = E \left[ \suminfn \frac{(\lambda x)^n}{n!} \right]
  = \suminfn \frac{\lambda^n E X^n}{n!}.
\]
It follows that $M_X' (0) = E X$. Now to justify the exchang of
integral and summation, note that
\[
  S_N = \sumi^N \frac{(\lambda x)^n}{n!} \to e^{\lambda x},
\]
and
\[
  \abs{S_N} \leq \sumi^N \frac{\abs{\lambda x}^n}{n!}
  \leq e^{\abs{\lambda x}}.
\]
However, $e^{\abs{\lambda X}} \leq e^{\lambda X} + e^{- \lambda X}$.
The expectation of RHS is finite for small $\lambda$ by assumption,
so the claim follows from dominated convergence theorem.

\begin{prop}
  Let $X$ be a non-negative random variable. Then TFAE:
  \begin{enumerate}
    \item $M_X$ is finite in some neighborhood of $0$.
    \item $\limsup_{p \to \infty} p^{-1} (E X^p)^{1 / p} < \infty$.
  \end{enumerate}
\end{prop}

\begin{proof}
  (1) $\implies$ (2). Suppose $M_X (\epsilon) = E \left[ \suminfn \frac{\epsilon^n
        X_n}{n!} \right] < \infty$. This implies that
  \[
    \sup_{n \geq 1} E \left[ \frac{\epsilon^n X^n}{n!} \right] < \infty.
  \]
  Let $1 \leq C < \infty$ be such that $E \left[ \frac{\epsilon^n X^n}{n!} \right] < C$
  for all $n \geq 1$. We then have
  \[
    \left( \frac{\epsilon}{n!} \right)^{1 / n} \left( E X^n  \right)^{1 / n}
    \leq C^{1 / n}
  \]
  It follows from Sterling's formula $n! \sim \left( n / e \right)^n$
  that
  \[
    \frac{1}{n} \left( E X^n \right)^{1 / n} \leq C
  \]
  for some other constant $C$.

  (2) $\implies$ (1). Suppose $\limsup_{p \to \infty} p^{-1} \left( E X^p \right)
    ^{1 / p} \leq C < \infty$. It follows that for any $n \geq 1$
  \[
    \frac{1}{n!} \left( E X^n \right) \leq C^n \frac{n^n}{n!} \leq C^n
  \]
  for some other constant $C$.
  Take $\epsilon = \frac{1}{2 C}$, we then have
  \[
    E \left[ \frac{\epsilon^n X^n}{n!} \right] \leq 2^{-n}.
  \]
  Therefore, $M_X (\epsilon)
    = E \left[ \suminfn \frac{\epsilon^n X^n}{n!} \right] < \infty$.

\end{proof}

Now we present an example of moment method for approximating random
variables.
\begin{eg}
  Let $G \sim G(n, \frac{1}{2})$ be the Erdos-Renyi random graph. The goal
  is to estimate the number of triangles in $G$. Let $N$ be
  the number of triangles in $G$. We have
  \[
    N = \sum_{\substack{S \subset [n] \\ \abs{S} = 3}} b_S,
  \]
  where $b_S$ is the indicator that $S$ is a triangle in $G$.
  It follows that $E N = \frac{1}{8} \binom{n}{3}$. Also,
  \[
    E N^2 = \sum_{\abs{S} = 3} \sum_{\abs{S'} = 3} E [b_S b_{S'}]
  \]
  To compute this, we consider several cases.
  \begin{enumerate}
    \item If $S \cap S' = \emptyset$, then $E [b_S b_{S'}] = \frac{1}{64}$.
    \item If $\abs{S \cap S'} = 1$, then $E [b_S b_{S'}] = \frac{1}{64}$.
    \item If $\abs{S \cap S'} = 3$, then $E [b_S b_{S'}] = \frac{1}{8}$.
    \item If $\abs{S \cap S'} = 2$, then $E [b_S b_{S'}] = \frac{1}{32}$.
  \end{enumerate}
  Hence,
  \[
    E N^2 = \frac{1}{64} \binom{n}{3}^2 \pm O(n^4),
  \]
  where the $O(n^4)$ term comes from the cases where $\abs{S \cap S'} = 2$
  or $3$. Therefore,
  \[
    \var N = E N^2 - (E N)^2 = O(n^4).
  \]
  Using Chebyshev's inequality, for each $t > 0$ we have
  \[
    P \left( \abs{N - E N} \geq t \cdot E N \right)
    \leq \frac{\var N}{t^2 (E N)^2} = t^{-2} O(n^{-2})
    \to 0.
  \]
  as $n \to \infty$.

\end{eg}

We have the following proposition on subgaussian decay of moments.
\begin{prop}
  Let $X$ be a random variable, TFAE:
  \begin{enumerate}
    \item $\sup_{p \geq 1} (E \abs{X}^p)^{1 / p} p^{-1/2} < \infty$.
    \item there exists $c > 0$ such that $P \left( \abs{X} > t \right)
            \leq 2 \exp (-c t^2)$ for all $t > 0$.
  \end{enumerate}
  If any of these statement is satisfied, $X$ is called a \emph{subgaussian
    variable}.
\end{prop}

\begin{proof}
  Exercise.

\end{proof}

\begin{thm}[Khintchine's inequality]
  There exists a universal constant $C < \infty$ such that for any
  $n \in \N$, $a_1, \dots, a_n \in \R$, and $p \geq 1$, we have
  \[
    \left( E \abs{\sumi^n a_i r_i}^p \right)^{1 / p} \leq
    C \sqrt{p} \norm{a}_2,
  \]
  where $r_1, \dots, r_n$ are i.i.d.\ Rademacher variables:
  $P \left( r_i = 1 \right) = P \left( r_i = -1 \right)
    = \frac{1}{2}$.
\end{thm}

\begin{proof}
  WLOG assume $\norm{a}_2= 1$. Let $\lambda > 0$. For any $t > 0$, we have
  \begin{align*}
    P \left( \sumi^n a_i r_i > t \right)
    = P \left( \exp \left( \lambda \sumi^n a_i r_i \right)
    > \exp (\lambda t) \right)
    \leq \frac{E \exp \left( \lambda \sumi^n a_i r_i \right)}
         {\exp (\lambda t)}
  \end{align*}
  from Markov's inequality. By independence, we have
  \begin{align*}
    P \left( \sumi^n a_i r_i > t \right)
    \leq \frac{\prod_{i=1}^{n} E \exp (\lambda a_i r_i) }{\exp (\lambda t)}
    = \frac{\prod_{i=1}^n \frac{1}{2} \left( \exp (\lambda a_i) +
        \exp (- \lambda a_i) \right) }{\exp (\lambda t)}.
  \end{align*}
  Now we use $\exp (x) \leq x + \exp (x^2)$ to obtain
  \begin{align*}
    P \left( \sumi^n a_i r_i > t \right)
    \leq \frac{\prod_{i=1}^n \exp (\lambda^2 a_i^2)}{\exp (\lambda t)}
    = \exp \left( \lambda^2 \norm{a}^2_2 - \lambda t \right).
  \end{align*}
  This holds for any $\lambda > 0$. Substituting $\lambda =
    \frac{1}{2} t \norm{a}^{-2}_2$, we have
  \[
    P \left( \sumi^n a_i r_i > t \right) \leq \exp \left( -\frac{1}{4}
    \frac{t^2}{\norm{a}^2_2} \right)
  \]
  Similarly we can bound $P \left( \sumi^n a_i r_i < -t \right)$
  as the Rademacher variables are symmetric. Hence we have
  \[
    P \left( \abs{\sumi^n a_i r_i} > t \right) \leq 2
    \exp \left( -\frac{1}{4} t^2 \right).
  \]
  The theorem follows from the previous proposition.

\end{proof}

\begin{thm}[Anti-concentration inequalities]
  Let $X$ be a non-negative random variable and $1 \leq p < q < \infty$.
  Suppose $\left( E X^p \right)^{1 / p} \leq C \left( E X^q \right)^{1 / q}$
  for some constant $C$. Then there exsits $\epsilon > 0$ depending
  only on $p$, $q$, $C$ such that
  \[
    P \left( X \geq \epsilon \left( E X^p \right)^{1 / p} \right)
    \geq \epsilon.
  \]
\end{thm}

\subsection{Convergence of random variable}

\begin{defi}
  Let $\seqinfn{X_n}$ be a sequence of random variables and $X$ is also
  a random variable.
  Say $X_n$ converges to $X$ \emph{in distribution} or \emph{weakly} if
  for any compactly supported continuous function $f$, we have
  \[
    \int_\R f(t) \d \mu_{X_n} (t) \to \int_\R f(t) \d \mu_X (t).
  \]
  Write $X_n \convd X$ or $X_n \convw X$.
\end{defi}

\begin{prop}
  $X_n \convw X$ iff $F_{X_n} (t) \to F_X (t)$ at every point $t$
  where $F_X$ is continuous.
\end{prop}

\begin{defi}
  Let $\seqinfn{X_n}$ be a sequence of random variables and $X$ is also
  a random variable defined on the same probability space.
  Say $X_n$ converges to $X$ \emph{in probability} if
  for any $\epsilon > 0$,
  \[
    P \left( \abs{X_n - X} < \epsilon \right) \to 1.
  \]
  Write $X_n \convp X$.
\end{defi}

\begin{prop}
  If $X_n \convp X$, then $X_n \convd X$.
\end{prop}

\begin{proof}
  Let $t \in \R$ be such that $F_X$ is continuous at point $t$
  and let $\epsilon > 0$. Let $\delta > 0$ be such taht
  $F_X (t + \delta) \leq F_X (t) + \epsilon$ and
  $F_X (t - \delta) \geq F_X (t) - \epsilon$.
  Since $P \left( \abs{X - X_n} < \delta\right)
    \to 1$, there exists $n_0$ such that $n \geq n_0$ implies
  \[
    P \left( \abs{X_n - X} < \delta\right)
    \geq 1 - \epsilon.
  \]
  Take $n \geq n_0$. We have
  \begin{align*}
    F_{X_n} (t)
     & = P \left( X_n \leq t \right)               \\
     & \leq P \left( \abs{X_n - X} > \delta\right)
    + P \left( X \leq t + \delta\right)            \\
     & \leq \epsilon + F_X (t) + \epsilon          \\
     & \leq F_X (t) + 2 \epsilon.
  \end{align*}
  Similarly for the other side.

\end{proof}

\begin{prop}
  If $X_n \to X$ a.s., then $X_n \convp X$.
\end{prop}

\begin{proof}
  Let $\epsilon > 0$. For each $\omega \in \Omega$, define
  \[
    n_\epsilon (\omega) = \min \left\{ n : \abs{X_m(\omega) - X(\omega)}
    \leq \epsilon \text{ for all $m \geq n$} \right\} < \infty.
  \]
  By continuity of measure, for any $\delta > 0$ there exists $N$ such that
  $P \left( n_\epsilon \leq N \right) \geq 1 - \delta$.
  This implies that for all $m \geq N$, we have
  \[
    P \left( \abs{X_m - X} \leq \epsilon \right) \leq 1 - \delta.
  \]
  Since $\delta > 0$ is arbitrary, this compeltes the proof.

\end{proof}

\begin{prop}
  Let $c \in \R$ and suppose $\seqinfn{X_n}$ are random variables
  on a common probability space
  $(\Omega, \Sigma, P)$. Then $X_n \convd c$ iff $X_n \convp c$.
\end{prop}

\begin{proof}
  Just need to check that $X_n \convd c$ implies $X_n \convp c$.
  We have
  \[
    F_c (t) = \begin{cases}
    1 & \text{if $t \geq c$,} \\
    0 & \text{otherwise.}
    \end{cases}
  \]
  Take $\epsilon > 0$, then $c \pm \epsilon$ are points of continuity.
  It follows that $F_{X_n} (c + \epsilon) \to 1$ and $F_{X_n} (c - \epsilon)
    \to 0$. This implies that
  \[
    P \left( \abs{X_n - c} > \epsilon \right) \to 0.
  \]

\end{proof}


Next we consider algebraic operations and convergence.
It is clear that if $X_n \to X$ and $Y_n \to Y$ a.e.\, then $X_n + Y_n
  \to X + Y$ a.e.\ Also, if $X_n \convp X$ and $Y_n \convp Y$, then
$X_n + Y_n \convp X + Y$. But for convergence in distribution,
the same statement need not be true without any extra assumptions.
However, suppose $X_n$ and $Y_n$ are independent for each $n$, $X$ and $Y$
are independent, $X_n \convd X$, and $Y_n \convd Y$. Then $X_n + Y_n
  \convd X + Y$.


\begin{eg}
  Suppose $\seqinfi{a_i}$ are i.i.d.\ random variables and
  \[
    \epsilon_n \sumi^n a_i \convd X
  \]
  for some random variable $X$ and sequence of positive
  numbers $\seqinfn{\epsilon_n}$.
  Then we have
  \[
    \epsilon_{2n} \sumi^{2n} a_i \convd X,
  \]
  and
  \[
    \frac{\epsilon_{2n}}{\epsilon_n} \left( \epsilon_n \sumi^n a_i \right)
    + \frac{\epsilon_{2n}}{\epsilon_n} \left( \epsilon_n \sum_{i=n+1}^{2n} a_i \right)
    \convd X.
  \]
  If we further suppose that there is $\alpha > 0$ such that
  \[
    \lim_{n \to \infty} \frac{\epsilon_{2n}}{\epsilon_n} = \alpha,
  \]
  we would have $\alpha X + \alpha \tilde{X} \sim X$, where
  $\tilde{X}$ is an independent copy of $X$. The only such distribution
  $X$ with bounded second moment is the Gaussian distribution.

\end{eg}

Recall that for independent random variables, we have a representation
theorem that create new independent random variables on the product space.
We alos have a representation theorem for random variables converging
in distribution. On the new probability space, they will coneverge a.e.\

\begin{thm}[Skorokhod's representation theorem]
  Suppose $X_n \convd X$ on some probability space $(\Omega, \Sigma, P)$.
  Then there exists another probability space $(\tilde{\Omega},
    \tilde{\Sigma}, \tilde{P})$ and random variables $\tilde{X}_n$,
  $\tilde{X}$ on it such that the following holds:
  \begin{itemize}
    \item $\tilde{X}_n \sim X_n$ for all $n$.
    \item $\tilde{X} \sim X$.
    \item $\tilde{X}_n \to \tilde{X}$ a.e.\
  \end{itemize}
\end{thm}

\begin{proof}
  Suppose for simplicity that $F_{X_n}$ and $F_X$ are all continuous and
  strictly increasing. With this, we can easily construct random
  variables with the same distribution. Let $U \sim \unif([0, 1])$,
  then $F_X^{-1} (U) \sim X$. Indeed,
  \[
    P \left( F_X^{-1} (U) \leq t \right)
    = P \left( F_X (F_X^{-1} (U)) \leq F_X (t) \right)
    = F_X (t).
  \]

  Take $([0, 1], \bcal_{[0, 1]}, m)$.
  Let $\tilde{X}_n = F_{X_n}^{-1} (\omega)$
  and $\tilde{X} = F_X^{-1} (\omega)$. Then
  $\tilde{X}_n \sim X_n$ and $\tilde{X} \sim X$.
  Note that $X_n \convd X$
  implies $F_{X_n} \to F_X$ pointwise. Checking
  $\tilde{X}_n \to \tilde{X}$ a.e.\ is left as an exercise.

\end{proof}

\subsection{Law of rare events}

As an example for convergence of random variables, we have the following
example presenting the law of rare events. First we need some definitions.

\begin{defi}[Poisson distribution]
  Say random variable $X \sim \Poisson (\lambda)$ has Possion distribution
  with parameter $\lambda$ for $\lambda > 0$ if
  \[
    P \left( X = m \right) = \frac{\lambda^m e^{- \lambda}}{m!}.
  \]
  Then, $E X = \lambda$, $E X^2 = \lambda^2 + \lambda$, and
  $\var X = \lambda$.
\end{defi}

\begin{defi}[exponential distribution]
  Say random variable $X \sim \Exp (\lambda)$ has exponential distribution
  with parameter $\lambda$ if $X$ has density function
  \[
    \rho_X (t) = \lambda \exp (- \lambda t) \quad \text{for $t \geq 0$.}
  \]
  Then the corresponding CDF is
  \[
    F_X (t) = \begin{cases}
    1 - \exp (- \lambda t), & \text{ if $t \geq 0$,} \\
    0, & \text{ otherwise.}
    \end{cases}
  \]
\end{defi}

The exponential distribution is a memoryless distribution: suppose
$X \sim \Exp (\lambda)$ and $t > s > 0$. Then,
\[
  P \left( X \geq t | X \geq s \right)
  = \frac{P \left( X \geq t \right)}{P \left( X \geq s \right)}
  = \exp (- \lambda (t - s))
  = P \left( X \geq t - s \right).
\]

\begin{eg}[law of rare events]
  Let $\lambda > 0$ be a fixed parameter. Let $\seqinfn{m_n}$
  be a non-decreasing sequence of integers such that
  \[
    \lim_{n \to \infty} \frac{m_n}{n} = \lambda.
  \]
  Create a triangular array as follows: for any $n$,
  let $\left\{ U_{n, j} \right\}_{j=1}^{m_n}$ be i.i.d.\ uniform random
  variables on $[0, n]$. Now let
  \[
    X_n = \abs{ \left\{ i \leq m_n : U_{n, i} \in [0, 1] \right\} }.
  \]
  Then $X_n \convd \Poisson (\lambda)$.

  To prove this, let $k \geq 0$, we have
  \begin{align*}
    P (X_n = k)
     & = \binom{m_n}{k} \left( \frac{1}{n} \right)^k
    \left( 1 - \frac{1}{n} \right)^{m_n - k}                                    \\
     & = (1 + o(1)) \frac{m_n^k}{k!} n^{-k} \exp \left( - \frac{m_n}{n} \right) \\
     & \to (1 + o(1)) \frac{\lambda^k e^{- \lambda}}{k!}
  \end{align*}
  as $n \to \infty$. Therefore, $X_n \convd \Poisson (\lambda)$ by the
  follwoing proposition.

\end{eg}

\begin{prop}
  Let $\seqinfn{X_n}$ and $X$ be integer-valued random variables.
  Then $X_n \convd X$ iff $f_{X_n} (k) \to f_X (k)$ for any $k \in \Z$.
\end{prop}

\begin{proof}
  Verify directly using compactly supported continuous test functions.
  Since the test functions
  are compactly supported, only finite sums are involved.

\end{proof}

\begin{eg}
  Continuing the previous example, define
  \[
    Y_n = \min_{1 \leq i \leq m_n} U_{n, i}.
  \]
  Then $Y_n \convd \Exp (\lambda)$.

  Let $t > 0$, then
  \begin{align*}
    P \left( Y_n \leq t \right)
     & = 1 - P \left( Y_n > t \right)                          \\
     & = 1 - P \left( \bigcap_{i=1}^{m_n} U_{n, i} > t \right) \\
     & = 1 - P \left( U_{n, 1} > t \right)^{m_n}               \\
     & = 1 - \left( \frac{n - t}{n} \right)^{m_n}              \\
     & \to 1 - \exp (- \lambda t)
  \end{align*}
  as $n \to \infty$.

\end{eg}

\begin{eg}[Poisson process]
  Fix integer parameter $w > 0$. For every large $n$, let
  \[
    Z_n = (U_{n, 1}^*, U_{n, 2}^*, \dots, U_{n, w}^*)
  \]
  be $w$ smallest numbers from the set $\left\{ U_{n, i} \right\}_{i=1}^{m_n}$,
  arranged in increasing order, then $Z_n$ is a random
  vector in $\R^w$. We want to investigate the limiting distribution
  of $\seqinfn{Z_n}$, if any exists.

  \textbf{Claim:} let $\xi_1$, \dots, $\xi_w \sim \Exp (\lambda)$ be i.i.d.\
  Let the random vector
  \[
    \xi = (\xi_1, \xi_1 + \xi_2, \dots, \xi_1 + \dots + \xi_w).
  \]
  Then $Z_n \convd \xi$.

\end{eg}

\begin{eg}[Poisson process]
  Let $\seqinfn{X_n}$ be a sequence of i.i.d.\
  random variables that follows the exponential distribution
  $\Exp (\lambda)$. Define Poisson process $\{S_n\}_{n=0}^\infty$ via
  \[
    S_n = \sumi^n X_i.
  \]
  This has
  two properties:
  \begin{enumerate}
    \item For each interval of length $1$ on $\R_+$, the number of
          $S_n$'s within this interval follows the Poission distribution
          with parameter $\lambda$.

    \item For any finite collection of disjoint intervals, the variables
          corresponding to the number of $S_n$'s within the intervals are
          mutually independent. This is because of the memoryless property
          of the exponential distribution.
  \end{enumerate}
\end{eg}

\begin{lemma}
  If $c_j \to 0$, $a_j \to \infty$ and $a_j c_j \to \lambda$,
  then $(1 + c_j)^{a_j} \to e^{\lambda}$.
\end{lemma}

\begin{thm}[Central limit theorem, basic form]
  Let $\seqinfn{b_n}$ be i.i.d.\ random variables that are
  $\Bernoulli (\frac{1}{2})$. Define $S_n = \sumi b_i$.
  Then
  \[
    \frac{S_n - \frac{n}{2}}{\sqrt{n} / 2} \convd \ncal (0, 1).
  \]
\end{thm}

\begin{proof}
  Suppose $n$ and $k$ be some integer, then
  \begin{align*}
    P \left( S_{2n} = n + k \right)
     & = \frac{(2n) !}{(n + k) ! (n - k)!} 2^{-2n}
  \end{align*}
  Using Sterling's formula $n! \sim \sqrt{2 \pi n} (\frac{n}{e})^n$,
  we have
  \begin{align*}
    \frac{(2n) !}{(n + k) ! (n - k)!}
     & \sim \frac{ (2n)^{2n} }{(n + k)^{n + k} (n - k)^{n - k}}
    \frac{\sqrt{2 \pi n}}{\sqrt{2 \pi (n + k)} \sqrt{2 \pi (n - k)}}
  \end{align*}
  Hence,
  \begin{align*}
    P \left( S_n = n + k \right)
     & \sim \left( 1 + \frac{k}{n} \right)^{- n - k - \frac{1}{2}}
    \left( 1 - \frac{k}{n} \right)^{-n + k - \frac{1}{2}}
    \left( \pi n \right)^{\frac{1}{2}}.
  \end{align*}
  We also have
  \[
    \left( 1 + \frac{k}{n} \right)^{- n - k}
    \left( 1 - \frac{k}{n} \right)^{-n + k}
    = \left( 1 - \frac{k^2}{n^2} \right)^{-n}
    \left( 1 + \frac{k}{n} \right)^{-k}
    \left( 1 - \frac{k}{n} \right)^{k}.
  \]
  Letting $k = x \sqrt{n / 2}$, we know
  \begin{align*}
    \left( 1 - \frac{k^2}{n^2} \right)^{-n}
    \left( 1 + \frac{k}{n} \right)^{-k}
    \left( 1 + \frac{k}{n} \right)^{k}
     & = \left( 1 - \frac{x^2}{2 n} \right)^{-n}
    \left( 1 + \frac{x}{\sqrt{2n}} \right)^{- x \sqrt{n/2}}
    \left( 1 - \frac{x}{\sqrt{2n}} \right)^{- x \sqrt{n/2}} \\
     & \to \exp \left( \frac{x^2}{2} \right)
    \exp \left( - \frac{x^2}{2} \right)
    \exp \left( - \frac{x^2}{2} \right)                     \\
     & = \exp \left( - \frac{x^2}{2} \right)
  \end{align*}
  by the lemma above.

  Therefore, if $2k / \sqrt{2n} \to x$, then
  \[
    P \left( S_{2n} = n + k \right)
    \sim \left( \pi n \right)^{- 1 / 2} e^{- x^2 / 2}.
  \]
  Now,
  \begin{align*}
    P \left( a \leq \frac{S_{2n} - n}{ \sqrt{n / 2} } \leq b \right)
     & = P \left( n + a \sqrt{n / 2} \leq S_{2n} \leq n + b \sqrt{n / 2} \right) \\
     & = \sum_{k \in \left[ a \sqrt{n / 2}, b \sqrt{n / 2} \right] \cap \Z}
    P \left( S_{2n} = n + k \right).
  \end{align*}
  Chaging variables $x = k \sqrt{n / 2}$, we then have
  \begin{align*}
    P \left( a \leq \frac{S_{2n} - n}{\sqrt{n / 2}} \leq b \right)
     & = \sum_{x \in [a, b] \cap \left( 2\Z / \sqrt{2n} \right)}
         (2 \pi)^{- 1 / 2} e^{-x^2 / 2} (2 / n)^{1 / 2}  \\
     & \approx \int_a^b \frac{1}{\sqrt{2 \pi}} e^{-x^2 / 2} \d x,
  \end{align*}
  using the definition of Riemann integral. The proof is
  now complete.

\end{proof}

\begin{defi}[Multivariate Gaussian]
  There are two equivalent definitions of multivariate Gaussian.
  \begin{enumerate}
    \item $X \in \R^n$ has Gaussian distribution if there is non-random
          $n$-by-$n$ non-random matrix and $\mu \in \R^n$ such that
          $X \sim M G + \mu$, where $G$ is the standard Gaussian vector in $\R^n$.
          That is, $G = (G_1, \dots, G_n)$ are random vector with
          i.i.d.\ $\ncal(0, 1)$ components.

    \item $X \in \R^n$ has Gaussian distribution if for all
          non-random $y \in \R^n$, $\bra X, y \ket$ is a scalar
          Gaussian random variables.
  \end{enumerate}
\end{defi}

\begin{proof}
  \TODO

\end{proof}

\section{Law of large numbers}

\subsection{Weak law of large numbers}

\begin{thm}[Weak law of large numbers]
  Suppose $\seqinfn{X_n}$ are mutually independent, identically
  distributed random variables. Suppose also $E \abs{X_1} < \infty$.
  Let $S_n = \suminfn X_i$, then $S_n / n \convp E X_1$.
\end{thm}

\begin{proof}
  We prove this using truncation method.
  Let $\epsilon > 0$ and choose truncation level $M (\epsilon) > 0$
  satisfying the following:
  \begin{align*}
    E \abs{X_1 \ind_{\abs{X_1} > M}} < \epsilon^2.
  \end{align*}
  By dominated convergence theorem, there exists such an $M$.
  Note that this also implies
  \[
    \abs{E X_1 - E \left( X_1 \ind_{\abs{X_1} \leq M} \right)} < \epsilon.
  \]
  For each $n$, we have
  \[
    S_n = \sumk^n X_k \ind_{\abs{X_k} \leq M} +
    \sumk^n X_k \ind_{\abs{X_k} > M}.
  \]
  Write
  \[
    T_n = \sumk^n X_k \ind_{\abs{X_k} \leq M} \quad
    \text{and} \quad
    R_n = \sumk^n X_k \ind_{\abs{X_k} > M}.
  \]
  For the first term, we can use Chebyshev's inequality to obtain
  \begin{align*}
    P \left( \abs{\frac{T_n}{n} -
               E [X_1 \ind_{\abs{X_1} \leq M}]} \geq \epsilon \right)
     & \leq \frac{1}{\epsilon^2} \var \left( \frac{T_n}{n} \right) \\
     & = \frac{1}{n \epsilon^2} \var (X_1 \ind_{\abs{X_1} \leq M}) \\
     & \leq \frac{4 M^2}{n \epsilon^2}
  \end{align*}
  Similarly, for the second term, we have
  \begin{align*}
    P \left( \abs{\frac{R_n}{n}} \geq \epsilon \right)
     & \leq \frac{1}{\epsilon} E \abs{\frac{R_n}{n}} \leq \epsilon.
  \end{align*}
  It follows that
  \[
    P \left( \abs{\frac{S_n}{n} - E X_1} \geq
    2 \epsilon + \epsilon^2 \right)
    \leq \frac{4 M^2}{n \epsilon^2} + \epsilon.
  \]
\end{proof}

\begin{remark}
  Note that the only step we use independence is when we calculate
  $\var (T_n / n)$. In fact, pairwise independence is sufficient
  for this step.
\end{remark}

\subsection{Borel-Cantelli lemmas}

\begin{thm}[Borel-Cantelli lemmas]
  Let $\seqinfn{A_n}$ be a sequence of events and
  \[
    P \left( A_n \text{ i.o.} \right)
    = P \left( \capinfk \bigcup_{n=k}^\infty A_n \right)
    = P \left( \limsup A_n \right)
  \]
  \begin{enumerate}
    \item If $\suminfn P (A_n) < \infty$, then
          $P \left( A_n \text{ i.o.} \right) = 0$.

    \item Suppose also $A_n$'s are mutually independent
          and $\suminfn P (A_n) = \infty$, then
          $P \left( A_n \text{ i.o.} \right) = 1$.
  \end{enumerate}
\end{thm}

\begin{proof}
  \begin{enumerate}
    \item We have
          \begin{align*}
            P \left( A_n \text{ i.o.} \right)
             & = P \left( \suminfn \ind_{A_n} = \infty \right).
          \end{align*}
          However, $E [\suminfn \ind_{A_n}] = \suminfn P (A_n) < \infty$,
          so $P (A_n \text{ i.o.}) = 0$.

    \item We have
          \begin{align*}
            P \left( A_n \text{ i.o.} \right)
            = P \left( \cupinfk \bigcap_{n = k}^\infty A_n^c \right)
            \leq \suminfk P \left( \bigcap_{n=k}^\infty A_n^c \right).
          \end{align*}
          Using independence and $1 - t \leq \exp (- t)$,
          we obtain
          \[
            P \left( \bigcap_{n=k}^M A_n \right)
            = \prod_{n=k}^M (1 - P (A_n))
            \leq \exp \left( - \sum_{n=k}^M P (A_n) \right) \to 0
          \]
          as $M \to \infty$. Hence, $P \left( \bigcup_{n=k}^\infty A_n \right)
            = 1$ for all $k$. Since $\bigcup_{n=k}^\infty A_n \downarrow \limsup A_n$,
          it follows from monotone continuity that $P (A_n \text{ i.o.}) = 1$.
  \end{enumerate}
\end{proof}

\begin{prop}
  Let $\seqinfn{b_n}$ be mutually independent Bernoulli variables
  with parameters $\seqinfn{p_n}$ with $0 < p_n < 1$.
  Then, $b_n \to 0$ a.e.\ iff $\suminfn p_n < \infty$.
\end{prop}

\begin{proof}
  For any $\omega$, $b_n (\omega) \to 0$ as $n \to \infty$ iff
  $b_n (\omega) = 1$ for at most finitely many $n$'s.
  Let $A_n = \left\{ b_n = 1 \right\}$. Then,
  $b_n \to 0$ a.e.\ iff for a.e.\ $\omega$ is contained
  in finitely many $A_n$'s. That is, $P (A_n \text{ i.o.}) = 0$.
  The proposition then follows from Borel-Cantelli lemma.

\end{proof}

\subsection{Strong law of large numbers}

\begin{thm}
  Let $\seqinfn{X_n}$ be mutually independent identically distributed
  random variables. Suppose $E \abs{X_1} < \infty$
  and let $S_n = \frac{1}{n} \sumk^n X_k$
  Then, $S_n \to E X_1$ a.e.\
\end{thm}

\begin{proof}
  See notes for proof by Etemadi.
\end{proof}

\begin{eg}[Coupon collector]
  Consider array of random variable $\seqinfn{\seqinfj{X_{j, n}}}$.
  For each $n$, the sequence $\seqinfj{X_{j, n}}$
  are i.i.d.\ uniform on $[n] = \left\{ 1, \dots, n \right\}$.
  Let
  \[
    T_n = \min \left\{ m \geq 1 : \cupi^m \{X_{i, n}\} = [n] \right\}.
  \]
  We want to show $T_n / (n \log n) \convp 1$.

  To show this, we use second moment method.
  For each $k \leq n$, let $T_{n, k}$ be the first time $k$ coupons
  are collected. We then have $T_{n, 0} = 0$ and
  \[
    T_n = \sumk^n (T_{n, k} - T_{n, k - 1}).
  \]
  Note that $T_{n, k} - T_{n, k - 1}$ follows a geometric distribution.
  Hence,
  \begin{align*}
     & E \left( T_{n, k} - T_{n, k - 1} \right) = \frac{n}{n - k + 1}, \\
     & \var \left( T_{n, k} - T_{n, k - 1} \right) =
    \frac{k - 1}{n} \frac{1}{\left( 1 - \frac{k - 1}{n} \right)^2}.
  \end{align*}
  It follows that
  \[
    E T_n = n \sumk^n \frac{1}{n - k + 1} \sim n \log n.
  \]
  Also,
  \begin{align*}
    \var T_n
     & = \sumk^n \var \left( T_{n, k} - T_{n, k - 1} \right)
    = \sumk^n \frac{k - 1}{n} \frac{1}{\left( 1 - \frac{k - 1}{n} \right)^2}
    = o(n^2 \log^2 n)
  \end{align*}
  It then follows from Chebyshev's inequality that
  $T_n / (n \log n) \convp 1$.

\end{eg}

\begin{eg}
  Let $v_0 \in \left\{ -1, 1 \right\}^n$ and define a sequence
  of random vectors $\seqinfi{X_i}$ inductively as follows:
  define $X_0 = v_0$, and $X_{i + 1}$ is a random vector supported
  on $\left\{ -1, 1 \right\}^n$ obtained from $X_i$ by selecting
  a random nuber $m \in [n]$ uniformly at random and flipping
  $m$-th component of $X_i$ with probability $\frac{1}{2}$.

  Now $X_k$ is a random vector for each $k$.
  Denote the uniform distribution on $\left\{ -1, 1 \right\}^n$
  as $P$, and we are interested in
  \[
    \norm{X_k - P}_{TV}
    = \inf_{\xi, \tilde{X}} P \left( \xi \neq \tilde{X} \right),
  \]
  where the inf is taken over all pairs of random vectors
  $\xi$ and $\tilde{X}$ such that $\tilde{X} \sim X_k$ and
  $\xi \sim P$.

  \textbf{Claim:} to make sure that $\norm{X_k - P}_{TV} = o(1)$,
  we must have $k = \Omega (n \log n)$.

  This is similar to the coupon collector problem. We need to have selected
  at least $n (1 - O( n^{-1/2} ))$ indicies for the distribution
  of $X_k$ to be roughly uniform. With a similar derivation
  to the coupon collector problem, this implies $k = \Omega (n \log n)$.


\end{eg}

\section{Central Limit theorem}

\subsection{Characteristic Function}

\begin{defi}[characteristic function]
  Let $X$ be a random vector in $\R^n$, we define the
  characteristic function (ch.f.) of $X : \R^n \to \C$ via
  \[
    \phi_X (t) = E e^{i \bra X, t \ket}
  \]
  Suppose $X$ has density $\rho_X$, then
  \[
    \phi_X (t) = \int_{\R^n} \exp \left( i \bra s, t \ket \right)
    \rho_X (s) \d s.
  \]
  This is essentially the Fourier transform of $\rho_X$.
\end{defi}

\begin{prop}
  Let $X$ be a random variable, its characteristic function
  $\phi_X$ is uniformly continuous.
\end{prop}

\begin{proof}
  We have
  \begin{align*}
    \abs{\phi_X (t + \delta) - \phi_X (t)}
     & = \abs{ E e^{i \bra X, t \ket} \left( e^{i \bra x, \delta \ket} - 1 \right) }
    \leq E \abs{ e^{i \bra X, \delta \ket} - 1 }.
  \end{align*}
  Note that $\abs{e^{i \bra X, \delta \ket} - 1} \leq 2$. The proposition
  then immediately follows from dominated convergence theorem.
\end{proof}

\begin{prop}
  If $X$ is a random variable with $E \abs{X}^k < \infty$
  for some $k \geq 1$. Then $\phi_X$ is $k$ times continuous
  differentiable and
  \[
    \phi^{(k)}(0) = i^k E X^k.
  \]
\end{prop}

\begin{prop}
  Suppose $X$ is random variable and $a, b \in \R$. Then
  \[
    \phi_{aX + b} (t) = \exp (itb) \phi_X (at).
  \]
\end{prop}

\begin{prop}
  If $X_1, X_2, \dots, X_n$ are mutually independent and
  $S = X_1 + \dots + X_n$, then
  \[
    \phi_S (t) = \prod_{k=1}^n \phi_{X_k} (t).
  \]
\end{prop}

\begin{prop}
  Suppose $X \sim \ncal (0, 1)$, then
  \[
    \phi_X (t) = \exp \left( - \frac{t^2}{2} \right).
  \]
\end{prop}

\begin{proof}
  Suppose $Y \sim \ncal (0, 1)$, and $X$, $Y$ are independent.
  It follows that
  \[
    \phi_{X + Y} (t) = \phi_X (t) \phi_Y (t) = \phi_X (t)^2.
  \]
  Meanwhile,
  \begin{align*}
    \phi_{X + Y} (t)
     & = \frac{1}{2 \pi} \int_{-\infty}^\infty \int_{-\infty}^\infty
    \exp \left( i (x + y) t \right)
    \exp \left( - \frac{x^2}{2} \right)
    \exp \left( - \frac{y^2}{2} \right)
    \d x \d y
  \end{align*}
  Using change of variable with $u = x + y$, $v = x - y$, we obtain
  \begin{align*}
    \phi_{X + Y} (t)
     & = \frac{1}{4 \pi} \int_{-\infty}^\infty \int_{-\infty}^\infty
    \exp \left( -\frac{1}{2} \left( \frac{u^2}{2}
    + \frac{v^2}{2} \right) + i u t \right)
    \d u \d v                                                        \\
     & = \frac{1}{\sqrt{4 \pi}} \int \exp
    \left( - \frac{1}{4} u^2 + i t u \right) \d u                    \\
     & = \phi_X (\sqrt{2} t).
  \end{align*}
  Now we have a function equation
  \[
    \phi_X (\sqrt{2} t) = \phi_X (t)^2.
  \]
  This implies that $\phi_X (t) = \exp (- t^2 / 2)$.

\end{proof}


Now we present inversion formula, which shows the important correspondence
between the ch.f.\ $\phi_X$ and the induced Borel measure $\mu_X$
of a random variable.

\begin{thm}[inversion formula]
  Let $X$ be a random variable. Let $\phi_X$ be the ch.f.\ and $\mu_X$
  be the Borel measure on $\R$ induced by $X$. Suppose interval $(a, b]$
  is such that $\mu_X (\{ a, b \}) = 0$. Then,
  \[
    \mu_X ((a, b]) =
    \frac{1}{2 \pi} \lim_{T \to \infty}
    \int_{-T}^T \left[ \frac{\exp (- i t a) - \exp (- i t b)}{i t} \phi
      (t) \right] \d t.
  \]

  This shows that there is a unique correspondence between $\mu_X$ and
  $\phi_X$.
\end{thm}

\begin{proof}
  Some application of Fubini-Tonelli theorem.
\end{proof}

\begin{remark}
  Using the formula, we can compute $\mu_X ((a, b])$ in terms of
  $\phi_X$ even if $\mu_X (\{a, b\}) \neq 0$.
  Indeed, there are at most countably many $x$ such that $\mu_X (\{x\})
    \neq 0$, so there exists a sequence $\epsilon_m$ such that
  $\epsilon_m \downarrow 0$, and $\mu_X (\{a + \epsilon_m, b + \epsilon_m\})
    = 0$. Then
  \[
    \mu_X ((a, b]) = \lim_{m \to \infty} \mu_X ((a_m, b_m]).
  \]
  It is important that $\epsilon_m \downarrow 0$ so that the
  approximation through monotone continuity is correct.
\end{remark}

The same is true in the $n$-dimensional setting. Now
\[
  \phi_X (t) = E \exp \left( i \bra t, X \ket \right).
\]
The induced Borel measure $\mu_X$ in $\R^n$ is uniquely determined by the
ch.f.\ $\phi_X$.

\begin{thm}
  If $X$ has distribution density $\rho_X$, it can be expressed
  as
  \[
    \rho_X (t) = \frac{1}{2\pi} \int_\R \exp (- i t s) \phi_X (s) \d s.
  \]
\end{thm}

\begin{thm}[Levy]
  Let $\seqinfk{X_k}$ be a sequence of random vectors in $\R^n$ and
  let $\seqinfk{\phi_k}$ be the corresponding ch.f. Then the following
  holds
  \begin{enumerate}
    \item Suppose $X_k \convd X$ and let $\phi$ be the ch.f.\ of $X$.
          Then $\phi_k (t) \to \phi(t)$ as $k \to \infty$ for all $t \in \R^n$.

    \item Suppose there exists $\phi$ that is continuous at zero and
          $\phi_k(t) \to \phi(t)$ for all $t \in \R^n$. Then, there exists $X$
          such that $\phi_X = \phi$ and $X_k \convd X$.
  \end{enumerate}
\end{thm}

\subsection{Central Limit Theorem}

With the results above, we can prove the central limit theorem.
\begin{thm}[central limit theorem]
  Let $\seqinfn{X_n}$  be a sequence of i.i.d.\ random variable
  such that $E X_1$ and $\var X_1$ are well-defined and finite.
  Let
  \[
    S_n = \frac{1}{\sqrt{n}} \sumk^n (X_k - E X_k).
  \]
  Then $S_n \convd \ncal (0, \var X_1)$.
\end{thm}

\begin{proof}
  WLOG assume $E X_1 = 0$ and $\var X_1 = 1$. Then
  \[
    \phi_{S_n} (t)
    = \prod_{k=1}^n \phi_{\frac{X_k}{\sqrt{n}}} (t)
    = \left( \phi_{X_1} \left( \frac{t}{\sqrt{n}} \right) \right)^n
  \]
  Now, using Taylor's formula, we have
  \[
    \phi (\epsilon) = 1 - \frac{\epsilon^2}{2} + o(\epsilon^2).
  \]
  Therefore,
  \begin{align*}
    \phi_{S_n} (t)
    = \left( 1 + \frac{t^2}{2 n} + o \left( \frac{t^2}{n} \right) \right)^n
    \to \exp \left( - \frac{t^2}{2} \right),
  \end{align*}
  as $n \to \infty$ for any $t \in \R$. This completes the proof.

\end{proof}

\begin{remark}
  Note that not only the partial sum $S_n$ converges to normal
  in distribution, but the sup norm of the CDF of $S_n$ and the
  CDF of normal distribution also converges to zero with a rate
  of $1 / \sqrt{n}$. This is the optimal rate and is achieved by
  Bernoulli $1/2$ variables.
\end{remark}

Now we present another characterization of central limit theorem
by Lindeberg. This characterization might provide more insight
and intuition compared to the computational approach we used above
with characteristic functions.

\begin{defi}[Lindeberg condition]
  Let $X_1, X_2, \dots$ be mutually independent random variables
  with zero mean. Suppose $\sigma_n^2 = \var X_n < \infty$ and
  $b_n^2 = \sumi^n
    \sigma_i^2$. Then $X_i$'s satisfy the \emph{Lindeberg condition}
  if for any $\epsilon > 0$,
  \[
    \lim_{n \to \infty} b_n^{-2}
    \sumi^n E \left[ X_i^2 \ind_{\{\abs{X_i} > \epsilon b_n\}} \right] = 0.
  \]
\end{defi}

\begin{thm}[Lindeberg CLT]
  Let $X_1, X_2, \dots$ be i.i.d.\ random variables that satisfy
  the Lindeberg condition. Define
  \[
    S_n = \frac{1}{b_n} \sumi^n X_i.
  \]
  Then $S_n \convd \ncal (0, 1)$.
\end{thm}

\begin{proof}
  Use the definition of weak convergence in terms of test functions.
  WLOG, we can work with functions $h \in C^\infty_c (\R)$, i.e.\ smooth
  and compactly supported functions. The goal is to show
  $E h(S_n) \to E h (\ncal(0, 1))$ as $n \to \infty$.

  We need to use Taylor's formula:
  \[
    h (x + u) = h(x) + u h' (x) + \frac{u^2}{2} h''(x) +
    O \left( \min \left\{ \abs{u}^2, \abs{u}^3 \right\} \right)
  \]
  Also, write
  \[
    S_n^{(j)} = \sumi^j Y_j + \sum_{i = j + 1}^n X_i,
  \]
  where $Y_i, X_i$ are mutually independent and $Y_i \sim \ncal(0, \sigma_i^2)$.
  Note that $b_n^{-1} S_n^{(n)} \sim \ncal(0, 1)$ and
  $S_n^{(0)} = \sumi^n X_i$, so the goal now is to show
  \[
    \abs{E h \left( b_n^{-1} S_n^{(n)} \right) - E h \left( b_n^{-1} S_n^{(0)} \right)}
    = o(1).
  \]
  It suffices to show that for all $1 \leq j \leq n$,
  \[
    \abs{E h \left( b_n^{-1} S_n^{(j)} \right) - E h \left( b_n^{-1} S_n^{(j - 1)} \right)}
    = o(b_n^{-2} \sigma_j^2)
  \]
  Now fix $1 \leq j \leq n$ and let
  \[
    Z = \sumi^{j - 1} Y_i + \sum_{i = j + 1}^n X_i.
  \]
  Then, using Taylor's formula, we obatin
  \begin{align*}
    h \left( b_n^{-1} S_n^{(j)} \right)
     & =  h \left( b_n^{-1} Z \right)
    + h' \left( b_n^{-1} Z \right) \left( b_n^{-1} S_n^{(j)} - b_n^{-1} Z \right)
    + \frac{1}{2} h'' \left( b_n^{-1} Z \right) \left( b_n^{-1} S_n^{(j)} - b_n^{-1} Z \right)^2 \\
     & \quad + O \left( \min \left\{ b_n^{-2} (S_n^{(j)} - Z)^2,
    b_n^{-3} (S_n^{(j)} - Z)^3 \right\} \right).
  \end{align*}
  Note that $S_n^{(j)} - Z = Y_j$ is independent of $Z$.
  It follows that
  \begin{align*}
    E h \left( b_n^{-1} S_n^{(j)} \right)
     & = E h \left( b_n^{-1} Z \right)
    + \frac{1}{2} \frac{\sigma_j^2}{b_n^2} E h'' \left( b_n^{-1} Z \right)
    + O \left( E \min \left\{ b_n^{-2} \abs{Y_j}^{2},
    b_n^{-3} \abs{Y_j}^3 \right\} \right)
  \end{align*}
  Similarly,
  \begin{align*}
    E h \left( b_n^{-1} S_n^{(j - 1)} \right)
     & = E h \left( b_n^{-1} Z \right)
    + \frac{1}{2} \frac{\sigma_j^2}{b_n^2} E h'' \left( b_n^{-1} Z \right)
    + O \left( E \min \left\{ b_n^{-2} \abs{X_j}^{2},
    b_n^{-3} \abs{X_j}^3 \right\} \right)
  \end{align*}
  Hence,
  \[
    \abs{E h\left( b_n^{-1} S_n^{(j)} \right) - E h \left( b_n^{-1} S_n^{(j - 1)} \right)}
    = O \left( E b_n^{-3} \abs{Y_j}^3
    + E \min \left\{ b_n^{-2} \abs{X_j}^{2},
    b_n^{-3} \abs{X_j}^3 \right\} \right)
  \]

  For simplicity, consider the classical setting where $\sigma_j^2 = 1$
  for all $j$. Now $b_n = \sqrt{n}$. Note that the 3rd absolute moment
  of standard Gaussian is finite, we have
  \[
    E
    b_n^{-3} \abs{Y_j}^3
    \leq O(b_n^{-3}) = O(n^{-3/2}) = o(n^{-1}).
  \]
  For the second term, let $\epsilon > 0$ be given and consider truncation:
  \begin{align*}
    E \min \left\{ b_n^{-2} \abs{X_j}^2, b_n^{-3} \abs{X_j}^{3} \right\}
     & = E \min \left\{ b_n^{-2} \abs{X_j}^2, b_n^{-3} \abs{X_j}^{3} \right\}
    \left( \ind_{(\abs{X_j} > M)} + \ind_{(\abs{X_j} \leq M)} \right)         \\
     & \leq E \left( b_n^{-2} \abs{X_j}^2 \ind_{(\abs{X_j} > M)} \right)
    + E \left( b_n^{-3} \abs{X_j}^3 \ind_{(\abs{X_j} \leq M)} \right)
  \end{align*}
  However,
  \[
    E b_n^{-2} \abs{X_j}^2 \ind_{(\abs{X_j} > M)}
    = \frac{1}{n} E \left( X_j \ind_{\abs{X_j} > M} \right)^2 \to 0
  \]
  as $M \to \infty$. We can then choose $M$ such that
  \[
    E \left( b_n^{-2} X_j^2 \ind_{(\abs{X_j} > M)} \right) \leq \frac{\epsilon}{n}
  \]
  Meanwhile,
  \[
    E \left( b_n^{-3} \abs{X_j}^3 \ind_{(\abs{X_j} \leq M)} \right)
    \leq \frac{M^3}{n^{3/2}}.
  \]
  This shows that
  \[
    E \min \left\{ b_n^{-2} \abs{X_j}^2, b_n^{-3} \abs{X_j}^{3} \right\}
    = o(n^{-1})
  \]
  and the proof is complete.
\end{proof}

\section{Martingales}

\subsection{Information and independence}

\begin{prop}
  Suppose $\{\fcal_i\}_{i \in I}$ is a family of $\sigma$-fields,
  and $\fcal_i \subset \Sigma$ for some $(\Omega, \Sigma, P)$.
  Let $\ecal \in \Sigma$. Suppose
  for any finite $J \subset I$, $\ecal$ is independent from
  $\sigma (\{\fcal_j : j \in J\})$. Then $\ecal$ is independent from
  $\sigma (\{\fcal_i: i \in I\})$.
\end{prop}

\begin{proof}
  Let $\Lambda$ be the collection of events in $\Sigma$ that are independent
  from $\ecal$. Then $\Lambda$ is a $\lambda$-system from $\sigma$-addivity
  of $P$. Now let
  \[
    \Pi = \left\{ E \in \Sigma : E \in \sigma (\{ \fcal_j : j \in J \})
    \text{ for some $j \in J$ finite} \right\}.
  \]
  Then, $\Pi \subset \Lambda$. Moreover, if $E_1, E_2 \in \Pi$
  and let $J_1, J_2$ be the corresponding finite index set, then
  $E_1 \cap E_2 \in \sigma (\{ \fcal_j : j \in J_1 \cup J_2 \})$.
  Hence $E_1 \cap E_2 \in \Pi$ and $\Pi$ is a $\pi$-system.
  Therefore, by $\pi$-$\lambda$ theorem, we can conclude that
  $\Lambda \supset \sigma (\Pi) = \sigma (\{ \fcal_i : i \in I \})$.

\end{proof}

\begin{thm}[Kolmogorov's $0$-$1$ law]
  Let $(\Omega, \Sigma, P)$ be probability space.
  Suppose $\seqinfi{\fcal_i}$ be a sequence of mutually independent
  $\sigma$-field of $\Sigma$.
  Define tail $\sigma$-field:
  \[
    \tcal = \capinfk \sigma \left( \bigcup_{i=k}^\infty \fcal_i \right).
  \]
  Then $P(E) \in \{0, 1\}$ for any $E \in \tcal$.
\end{thm}

\begin{proof}
  For any $k \geq 1$, we have $E \in \sigma (\bigcup_{i=k}^\infty \fcal_i)$,
  so $E$ is independent from $\sigma (\fcal_1, \dots, \fcal_{k - 1})$.
  The verification of the previous claim is left as an exercise.
  Now from the previous proposition,
  this implies that $E$ is independent from $\sigma (\cupinfi \fcal_i)
    \supset \tcal$. Therefore, $E$ is independent from $\tcal$
  and thus $E$ is independent from itself. That is,
  \[
    P(E) = P(E)^2,
  \]
  and the desired results immediately follows.

\end{proof}

\begin{ex}
  Verify the claim in the proof above.
\end{ex}

\begin{proof}
  We just need to verify that if $\fcal_1, \dots, \fcal_n$ are
  independent, then $\fcal_1, \dots, \fcal_{k - 1},
    \sigma (\bigcup_{i = k}^n \fcal_i)$ are independent.
  Fix $E_1 \in \fcal_1, \dots, E_{k - 1} \in \fcal_{k - 1}$
  and let $\Lambda$ be the collection of sets $E \in \Sigma$ such
  that $P (E \cap E_1 \cap \dots \cap E_{k - 1})
    = P(E) P (E_1) \dots P(E_{k - 1})$.
  It can be easily verified that this is a $\lambda$-system.
  Now let $\Pi$ be the collection of sets in the form
  of $E_k \cap \dots \cap E_n$ where
  $E_k \in \fcal_k, \dots, E_n \in \fcal_n$.
  This is a $\pi$-system and generates $\sigma (\bigcup_{i=k}^n
    \fcal_i)$. Moreover, $\Pi \subset \Lambda$.
  The desired statement then follows from the $\pi$-$\lambda$ system.

\end{proof}

Now we present a simple applications of the Kolmogorov's $0$-$1$ law.

\begin{prop}
  Suppose $X_1, X_2, \dots$ are mutually independent random
  variables. Let $S_n = \sumi^n X_i$ and
  $C = \{ \lim_{n \to \infty} S_n / n \text{ exists} \}$.
  Then, $P (C) \in \{ 0, 1 \}$.
\end{prop}

\begin{proof}
  Fix $k \geq 1$ and let $n \geq k$, then
  \[
    \frac{S_n}{n}
    = \frac{X_1 + \dots + X_k}{n} + \frac{1}{n} \sum_{i = k + 1}^n X_i.
  \]
  The first term goes to zero as $n \to \infty$.
  Hence, $\lim_{n \to \infty} S_n / n$ exists
  iff $\lim_{n \to \infty} n^{-1} \sum_{i = k + 1}^n X_i$ exists.
  This implies that $C \in \sigma (\bigcup_{i = k + 1}^n X_i)$.
  Since this is true for all $k \geq 1$, $C$ is in the tail
  $\sigma$-field. The proposition then follows from Kolmogorov's
  $0$-$1$ law.

\end{proof}

\subsection{Conditional expectation}

The notion of conditional expectation is used to study independence
between current element of a random process and its history.

As a elementary definition and construction,
let $(\Omega, \Sigma, P)$ be a probability space.
Suppose $\{ A_j \}_{j \in J}$ be a finite partition of $\Omega$
into disjoint events that has positive probability.
Suppose also $X$ a random variable with $E \abs{X} <
  \infty$. For each $j$, we can consider the conditional probability
space and let $E_j$ be the expectation of the restriction of $X$.
Then,
\[
  E_j = \frac{1}{P (A_j)} E [X \ind_{A_j}]
\]
Now let $Y$ be a new random variable on $(\Omega, \Sigma, P)$ defined
via $Y (\omega) = E_j$ if $\omega \in A_j$. It follows that
\[
  E Y
  = \sum_{j \in J} E_j P(A_j)
  = \sum_{j \in J} E [X \ind_{A_j}]
  = E X.
\]
Moreover, let $I \subset J$ and let $A_I = \bigcup_{j \in I} A_j$,
then $\ind_{A_I} = \sum_{j \in I} \ind_{A_j}$.
We can similarly obtain $E [Y \ind_{A_I}] = E [X \ind_{A_I}]$.

Now we present the formal definition.
\begin{defi}[Conditional expectation]
  Let $(\Omega, \Sigma, P)$ be a probability space.
  Suppose $\fcal \subset \Sigma$ be $\sigma$-field and random
  variable $X$ with $E \abs{X} < \infty$.
  Then the conditional expectation of $X$ given $\fcal$
  is a $\fcal$-measurable random variable, write $E [X | \fcal]$
  such that for any $A \in \fcal$,
  \[
    \int_A X \d P = \int_A E [X | \fcal] \d P.
  \]
\end{defi}

\begin{remark}
  Note that if $\fcal$ is just $\{ \emptyset, \Omega \}$, then
  the conditional expectation $E [X | \fcal] = E X$.
\end{remark}

\begin{thm}
  Conditional expectation exists and is unique up to a null set.
\end{thm}

\begin{proof}
  (Uniqueness) Suppose both $f$ and $g$ are
  conditional expectation of $X$. Then
  \[
    \int_A f \d P = \int_A X \d P = \int_A g \d P
  \]
  for all $A \in \fcal$. Hence, $f = g$.

  (Existence) First suppose $X$ is nonnegative.
  Define measure $\nu$ on $\fcal$ via
  \[
    \nu (A) = \int_A X \d P
  \]
  for all $A \in \fcal$. This is a finite measure on $\fcal$,
  and it is absolutely continuous w.r.t.\ $P$.
  By Random-Nikodym theorem, there exists $\fcal$-measurable
  function $f \geq 0$ such that $\nu (A) = \int_A f \d P$.
  Therefore, $f$ is a conditional expectation of $X$ given $\fcal$.
  For real-valued $X$, let $f^\pm$ be the conditional expectation
  of $X^\pm$. Then, $f^+ - f^-$ is a conditional expectation for $X$.

\end{proof}

\begin{prop}
  We have the following properties of conditional expectation:
  \begin{enumerate}
    \item Suppose $E \abs{X} < \infty$, $E \abs{Y} < \infty$,
          and $X \leq Y$ a.e. Then, $E[X | \fcal] \leq E[Y | \fcal]$.

    \item Suppose $X_n \uparrow X$ a.e., $X \geq 0$ and
          $E \abs{X} < \infty$. Then, $E [X_n | \fcal] \uparrow
            E [X | \fcal]$ a.e.

    \item If $X$ is $\fcal$-measurable, then $E [X | \fcal]
            = X$.

    \item If $X$ is independent of $\fcal$, then
          $E [X | \fcal] = E X$.

    \item If $\fcal_1 \subset \fcal_2 \subset \fcal$, then
          $E[ E[X | \fcal_2] | \fcal_1] = E [X | \fcal_1]$.

  \end{enumerate}
\end{prop}

\begin{eg}
  Let $\seqinfi{X_i}$ be i.i.d.\ $\pm 1$ symmetric variable,
  then $\sumi^n X_i$ is the simple symmetric random walk on $\Z$.
  Define $\fcal_n = \sigma (X_1, \dots, X_n)$, then
  $\seqinfn{\fcal_n}$ forms a filtration. We hope to compute
  the conditional expectation $E [S_n | \fcal_{n - 1}]$.

  Since $S_n = X_n + S_{n - 1}$, we have
  \begin{align*}
    E[S_n | \fcal_{n - 1}]
     & = E[S_{n - 1} | \fcal_{n - 1}] + E[X_n | \fcal_{n - 1}] \\
     & = S_{n - 1} + E X_n                                     \\
     & = S_{n - 1}.
  \end{align*}

\end{eg}

\begin{remark}
  If $\seqinfn{Y_n}$ is a process such that $E (Y_n | Y_1, \dots,
    Y_{n - 1}) = Y_{n - 1}$, it is a discrete time martingale.
  In the example above $\seqinfn{S_n}$ is a discrete time martingale.
\end{remark}

\begin{remark}
  If $\seqinfn{Y_n}$ is on a countable state space and
  $F_{(Y_n | Y_1, \dots Y_{n - 1})} = F_{(Y_n | Y_{n - 1})}$,
  we say $\seqinfn{Y_n}$ satisfies Markov property and is a Markov process.
  In the example above $\seqinfn{S_n}$ is a Markov process.
\end{remark}

\begin{defi}[conditional probability]
  The conditional probability $P(A | \fcal)$ is defined via
  \[
    P (A | \fcal) = E [\ind_A | \fcal].
  \]
\end{defi}

\begin{prop}
  If $X$ and $Y$ are independent random variables, then
  $E[X | Y] = E X$.
\end{prop}

\begin{proof}
  For any event $A \in \sigma (Y)$, we have
  \begin{align*}
    \int_A E[X | Y] \d P = \int_A X \d P
    = \int X \ind_A \d P.
  \end{align*}
  However,
  \[
    \int X \ind_A \d P = E[X] E[\ind_A] = P(A) E[X] = \int_A E X \d P.
  \]
  By uniqueness of conditional expectation, the proof is complete.

\end{proof}

\begin{prop}
  If $X$ and $Y$ are independent random variables
  and $f$ is a Borel functions on $\R^2$,
  then
  \[
    E [f(X, Y) | X] = g(X),
  \]
  where $g(x) = E f(x, Y)$ for all $x \in \R$.
\end{prop}

\begin{proof}
  Consider a setting where $X(\omega_1, \omega_2) = h(\omega_1)$,
  $Y(\omega_1, \omega_2) = z(\omega_2)$ and $\Omega = \Omega_X \times
    \Omega_Y$. Then for any event $A \in \Omega_X$, we have
  \begin{align*}
    \int_A g(X)
     & = \int_A \int_{\Omega_Y}
    f(X(\omega_1), Y(\omega_2))  \d P(\omega_1) \d P(\omega_2) \\
     & = \int_{A \times \Omega_Y} f(X(\omega_1), Y(\omega_2))
    \d P(\omega_1, \omega_2).
  \end{align*}
\end{proof}

\subsection{Reflection principle}
Consider a random walk $S_n$ on $\Z$.
Let $\ell \leq 0$ and $u > \ell$ be integers and consider a $n$-time horizon.
Start with a realization of
$\seqinfk{S_k}$ such that $S_n = u$ and $S_k \leq k$ for some
$k \leq n$. Then there is a bijection between paths that end at $u$
and paths that end at $2 \ell - u$. Let $M_n = \min_{0 \leq k \leq n} S_k$
be the running minimum of the random walk, then we have
\[
  P \left( S_n = u, M_n \leq \ell \right)
  = P (S_n = 2l - u).
\]

\begin{eg}
  Now we can compute the distribution of $M_n$.
  Take $\ell \leq 0$, we have
  \begin{align*}
    P \left( M_n \leq \ell \right)
     & = P \left( M_n \leq \ell, S_n \leq \ell \right)
    + P \left( M_n \leq \ell , S_n > \ell \right)                              \\
     & = P \left( S_n \leq \ell \right)
    + P \left( M_n \leq \ell , S_n > \ell \right)                              \\
     & = P ( S_n \leq \ell )
    + \sum_{u = \ell + 1}^\infty P \left( M_n \leq \ell , S_n = u \right)      \\
     & = P ( S_n \leq \ell ) + \sum_{u = \ell + 1}^\infty P (S_n = 2 \ell - u) \\
     & = 2 P ( S_n < \ell ) + P (S_n = \ell).
  \end{align*}
  Therefore,
  \[
    P \left( M_n = \ell \right) = P(S_n = \ell) + P(S_n = \ell - 1).
  \]

  Then we can compute expectation:
  \begin{align*}
    E M_n
     & = \sum_{\ell = - \infty}^0 \ell P \left( M_n = \ell \right) \\
     & = \sum_{\ell = - \infty}^0 \ell P \left( S_n = \ell \right)
    + \sum_{\ell = - \infty}^{-1} (\ell + 1) P(S_n = \ell)         \\
     & = \sum_{\ell = - \infty}^{0} \ell P ( S_n = \ell )
    + \sum_{\ell = 1}^\infty (- \ell) P( S_n = \ell )
    + \sum_{\ell = - \infty}^{-1} P(S_n = \ell)                    \\
     & = E [- \abs{S_n}] + P (S_n < 0).
  \end{align*}
  As $n \to \infty$, we have
  \[
    \frac{E M_n}{E \abs{S_n}} \to -1.
  \]
  Hence,
  \[
    E M_n = - \sqrt{2 \pi^{-1}} \sqrt{n} + o( \sqrt{n} )
  \]
\end{eg}

\begin{ex}
  Use reflection principle to show that the distribution
  of the first return time $T = \min \{ n > 0: S_n = 0 \}$ is
  \[
    P(T = 2k) = \frac{1}{2 k - 1} 2^{-2 k} \binom{2k}{k}.
  \]
  Then show $P (T = \infty) = 0$. This is called the \emph{recurrence}
  property.
\end{ex}

\begin{eg}[recurrence]
  From the previous exercise, we know the first return time $T$
  is finite with probability one. It follows that
  \[
    P \left( S_n = 0 \text{ i.o.} \right) = 1.
  \]
  Let $T_1$ be the time of first return to zero, i.e.\
  $T_1 = \min \{ t \geq 1 : S_t = 0 \}$, and define
  a random sequence $\seqinfn{S_{T_1 + n}}$.
  Let $i_1, \dots, i_k$ be fixed integers, we want to show
  \[
    P \left( \cap_{j=1}^k S_{T_1 + j} = i_j \right)
    = P \left( \cap_{j=1}^k S_j = i_j \right).
  \]
  Towards using Markov property of random walk, we have
  \begin{align*}
    P \left( \cap_{j=1}^k S_{T_1 + j} = i_j \right)
     & = \sum_{\ell \geq 1}  \sum_{u_1, \dots, u_{\ell - 1}}
    P \left( T_1 = \ell, \cap_{j=1}^{\ell - 1} S_j = u_j,
    \cap_{j=1}^k S_{\ell+ j} = i_j \right)
  \end{align*}
  For each term, we have
  \begin{align*}
    P \left( T_1 = \ell, \cap_{j=1}^{\ell - 1} S_j = u_j,
    \cap_{j=1}^k S_{\ell + j} = i_j \right)
    = P \left( T_1 = \ell, \cap_{j=1}^{\ell - 1} S_j = u_j \right)
    P \left( \cap_{j=1}^k S_{\ell + j} = i_j | T_1 = \ell \right)
  \end{align*}
  The second term is the same for any $\ell \geq 0$, so the summation
  becomes
  \begin{align*}
    P(\cap_{j=1}^k S_{T_1 + j} = i_j)
     & = \sum_{\ell \geq 1} P(T_1 = \ell)
    P \left( \cap_{j=1}^k S_{\ell + j} = i_j | T_1 = \ell \right) \\
     & = P(\cap_{j=1}^k S_{T_1 + j} = i_j)
  \end{align*}
  Therefore, $S_{T_1 + n} \sim S_n$ are equidistributed.

  Following this, we can also define $T_k$ inductively
  as the time of $k$-th return to zero, and it is easy to
  see that $P (T_k < \infty \text{ for all $k$}) = 1$. Also, note that
  $\{ T_k \leq \ell \} \in \sigma (S_1, \dots, S_\ell)$, this is
  called the \emph{stopping time} property.
\end{eg}

\subsection{Subgaussian random variables}

The convariance matrix of $X \in \R^n$ with $E X = 0$ is defined as
\[
  \cov [X] = E [X X^T].
\]
Note that $\tr E[X X^T] =
  E [\tr X X^T] = E [X^T X] = E \norm{X}^2$.

\begin{defi}
  Let $X$ be a random vector in $\R^n$. Suppose $EX = 0$
  and $\cov [X] = \Sigma \in \R^{n \times n}$. If
  for any $y \in \R^n$ and $t > 0$,
  \[
    P \left\{ \bra X, y \ket > t ( E \bra X, y \ket^2 )^{1/2}
    \right\} \leq 2 \exp \left( - \frac{t^2}{k^2} \right)
  \]
  We say $X$ is a subgaussian vector.
\end{defi}

If $X$ is a subgaussian vector, then we have some constraints on
the norm of $X$.

\begin{thm}
  Suppose $X$ is a subgaussian random vector, then there exists
  $C_k, c_k > 0$ such that
  \[
    P \left\{ \norm{X}_2 > t ( E \norm{X}_2^2 )^{1/2} \right\}
    \leq C_k \exp (- c_k t^2)
  \]
  for all $t > 0$.
\end{thm}

\begin{proof}
  The idea of the proof is to use test vectors and conditional
  expectation. Let $V \sim \ncal (0, I_n)$ and independent from $X$.
  Fix $t > 0$ and define event
  $A_t = \{ \norm{X} > t (E \norm{X}^2)^{1/2} \}$.
  There exists some $c$ such that for any $\tilde{x} \in \R^n$,
  $P \{ \abs{\bra V, \tilde{x} \ket} \geq c \norm{\tilde{x}} \} \geq
    \frac{3}{4}$. Define function $f$ and $g$ via
  \[
    f(\tilde{x}, \tilde{v}) = \ind_{\left\{ \abs{\bra \tilde{v},
        \tilde{x} \ket} \geq c \norm{\tilde{x}} \right\}},
    \quad
    g(\tilde{x}) = E f(\tilde{x}, V)
    = P \{ \abs{\bra V, \tilde{x} \ket} \geq c \norm{\tilde{x}}\}
    \geq \frac{3}{4}.
  \]
  Since $V$ and $X$ are independent, $E [f(X, V) | X] = g(X)$ a.e.
  It follows that
  \[
    P \{ \abs{\bra V, X \ket} \geq c \norm{X} | X \} \geq \frac{3}{4} \text{ a.e.}
  \]
  Next, since the componenets of $V$ are i.i.d.\ $\ncal(0, 1)$,
  we have $E (V^T \Sigma V) = \tr \Sigma = E \norm{X}^2$.
  Note that $\Sigma$ is PSD, Markov inequality gives
  \[
    P \{ V^T \Sigma V \leq 4 E \norm{X}^2 \}
    = P \{ V^T \Sigma V \leq 4 E \norm{X}^2 | X \}
    \geq \frac{3}{4}.
  \]
  Combining the two inequalities, we obtain
  \[
    P \{ V^T \Sigma V \leq 4 E \norm{X}^2 \text{ and }
    \abs{\bra V, X \ket} \geq c \norm{X} | X \}
    \geq \frac{1}{2} \text{ a.e.}
  \]
  Define event $A = \{ V^T \Sigma V \leq 4 E \norm{X}^2 \text{ and }
    \abs{\bra V, X \ket} \geq c \norm{X} \}$.

  Now we have $P(A \cap A_t) = P(A_t) P(A | A_t)$. Note that
  $A_t$ is $X$-measurable and $P( A | X ) \geq \frac{1}{2}$
  a.e. Therefore, $P(A \cap A_t) \geq \frac{1}{2} P(A_t)$.
  Through computation we can obtain
  \[
    A \cap A_t \subset \left\{ \abs{ \bra V, X \ket }
    \geq \frac{1}{2} ct \sqrt{V^T \Sigma V} \right\}.
  \]
  We now just need to bound the probability of the event on RHS.

  For any $v \in \R^n$, define a new function $h$ via
  \begin{align*}
    h (v)
    = E \left[ \ind_{ \left\{ \abs{\bra v, X \ket} \geq \frac{1}{2}
            c t \sqrt{ v^T \Sigma v} \right\} } \right]
    = E \left[ \ind_{ \left\{ \abs{\bra v, X \ket} \geq \frac{1}{2} c t
            \sqrt{E \bra v, X \ket^2} \right\} } \right]
  \end{align*}
  Since $X$ is subgaussian, $h(v) \leq 2 \exp \left( - \frac{c^2 t^2}{4 k^2} \right)$ a.e.
  Therefore,
  \[
    E \left[ P \left\{ \abs{V, X} \geq \frac{1}{2} c t \sqrt{V^T \Sigma V} \middle| V \right\} \right]
    = E h (V) \leq 2 \exp \left( - \frac{c^2 t^2}{4 k^2} \right).
  \]
  The LHS is exactly $P \{ \abs{\bra V, X \ket} \geq \frac{1}{2}
    c t \sqrt{ V^T \Sigma V } \} $, so
  \[
    P(A_t) \leq 2 P(A \cap A_t) \leq 4 \exp \left( - \frac{c^2 t^2}{4 k^2} \right)
  \]

\end{proof}

\subsection{Martingales and concentration inequalities}

\begin{defi}
  Let $\seqinfn{\fcal_n}$ be a filtration and $\seqinfn{X_n}$ be
  an adapted sequence, i.e.\ $X_n$ is $\fcal_n$ measurable.
  Suppose also $E \abs{X_n} < \infty$ for all $n$. Then
  \begin{enumerate}
    \item $\seqinfn{X_n}$ is a submartingale if
          $E [X_{n + 1} | \fcal_n] \geq X_n$ for all $n$.
    \item $\seqinfn{X_n}$ is a supermartingale if
          $E [X_{n + 1} | \fcal_n] \leq X_n$ for all $n$.
    \item $\seqinfn{X_n}$ is a martingale if
          $E [X_{n + 1} | \fcal_n] = X_n$ for all $n$
  \end{enumerate}
\end{defi}

\begin{eg}
  Suppose $\seqinfn{Y_n}$ are mutually independent nonnegative
  random variables with $E Y_n = 1$. Define $X_n = \prod_{j=1}^n Y_j$.
  Then $\fcal_n = \sigma (Y_1, \dots, Y_n)$ is a filtration,
  and with this filtration $X_n$ is a martingale.
\end{eg}

\begin{eg}
  Suppose $X$ is random variable with $E \abs{X} < \infty$
  and $\fcal_0 \subset \fcal_1 \subset \cdots$ is a filtration.
  Then $X_n = E [X | \fcal_n]$ is a martingale.
\end{eg}

Next we are going to present some concentration inequalities
based on martingales.

\begin{thm}[Azuma's inequality]
  Let $(\Omega, \Sigma, P)$ be a probability space.
  Suppose $\seqi{\fcal_i}^N$ be a filtration with $N < \infty$.
  Suppose $\seqi{X_i}^N$ is a martingale
  and suppose the $L^\infty$ norm
  $d_i = \norm{X_i - X_{i - 1}}_\infty < \infty$ for any $i \leq N$.
  Then
  \[
    P \left\{ \abs{X_N - X_0} \geq t \right\}
    \leq 2 \exp \left( - \frac{t^2}{4 \sumi^N d_i^2} \right)
  \]
  for all $t > 0$.
\end{thm}

\begin{proof}
  We aim to estimate $P \left\{ X_N - X_0 \geq t \right\}$.
  Fix $\lambda > 0$ and we have
  \begin{align*}
    E \exp (\lambda (X_N - X_0))
     & = E \exp \left( \lambda \sumi^N (X_i - X_{i - 1}) \right)                      \\
     & = E \left[ E \left[ \exp \left( \lambda \sumi^N (X_i - X_{i - 1})
               \right) \middle|  \fcal_{N - 1} \right] \right]              \\
     & = E \left[ \exp \left( \lambda \sumi^N (X_i - X_{i - 1}) \right)
             E \left[ \exp (\lambda (X_N - X_{N - 1})) | \fcal_{N- 1} \right] \right]
  \end{align*}
  No we need to use martingale property of $X_n$. However, we
  are dealing with the expectation of an exponential. Using
  $\exp x \leq x + \exp (x^2)$ for all $x \in \R$, we have
  \begin{align*}
    E [\exp (\lambda (X_N - X_{N - 1})) | \fcal_{N - 1}]
     & \leq E [\lambda (X_N - X_{N - 1}) + \exp (\lambda^2 (X_N - X_{N - 1})^2)
                | \fcal_{N - 1}] \\
     & = E [\exp (\lambda^2 (X_N - X_{N - 1})^2) | \fcal_{N - 1}]               \\
     & \leq \exp (\lambda^2 d_N^2).
  \end{align*}
  Therefore,
  \begin{align*}
    E \exp (\lambda (X_N - X_0))
     & \leq \exp (\lambda^2 d_N^2)
    E \left[ \exp \left( \lambda \sumi^N (X_i - X_{i - 1}) \right) \right],
  \end{align*}
  and inductively we can obtain
  $E \exp (\lambda (X_N - X_0)) \leq \exp \left( \lambda^2 \sumi^N d_i^2 \right)$.
  Finally, Markov inequality gives
  \[
    P \left\{ \exp (\lambda (X_N - X_0)) \geq \exp (\lambda t) \right\}
    \leq \frac{E \exp (\lambda (X_N - X_0))}{\exp (\lambda t)}
    \leq \exp \left( \lambda^2 \sumi^N d_i^2 - \lambda t \right)
  \]
  Letting $\lambda = \frac{1}{2} t (\sumi^N d_i^2)^{-1}$ completes the
  proof.

\end{proof}

\begin{cor}[Bounded differences inequality]
  Let $\seqi{\fcal_i}^N$ be a filtration with $\fcal_N = \Sigma$ and
  $\fcal_0 = \{ \emptyset, \Omega \}$.
  Let $\xi$ be a random variable with $E \abs{\xi} < \infty$.
  Let $d_i = \norm{ E[\xi | \fcal_i] - E [\xi | \fcal_{i - 1}]}_\infty$
  for $1 \leq i \leq N$. Suppose that $d_i < \infty$ for all
  $i \leq N$, then
  \[
    P \left\{ \abs{\xi - E \xi} > t \right\}
    \leq 2 \exp \left( - \frac{t}{4 \sumi^N d_i^2} \right)
  \]
  for all $t > 0$.
\end{cor}

\begin{cor}
  Suppose $Y_1, \dots, Y_N$ are mutually independent random
  variables. Suppose $f : \R^N \to \R$ satisfies the
  following ``Lipschitz'' condition: for any
  $i \leq N$ and any choice of $x_j \in \ran (Y_i)$
  for $j \neq i$ and $x_i, \tilde{x_i} \in \ran (Y_i)$,
  \[
    \abs{f(x_1, \dots, x_{i - 1}, x_i,
      x_{i + 1}, \dots, x_N) - f(x_1, \dots, x_{i - 1}, \tilde{x_i},
      x_{i + 1}, \dots, x_N)} \leq 1.
  \]
  Then
  \[
    P \left\{ \abs{f(Y_1, \dots, Y_N) - E f(Y_1, \dots, Y_N)} > t \right\}
    \leq 2 \exp \left( - \frac{t^2}{4 N} \right)
  \]
  for any $t > 0$.
\end{cor}

\begin{proof}
  Consider a filtration given by $\fcal_0 = \{ \emptyset, \Omega \}$,
  $\fcal_N = \Sigma$, and $\fcal_i = \sigma(Y_1, \dots, Y_i)$
  for $1 \leq i < N$. Let $\xi = f (Y_1, \dots, Y_N)$ and consider
  \begin{align*}
    \abs{E [\xi | \fcal_i] - E [\xi | \fcal_{i - 1}]}
     & = \abs{E_{Y_{i+1}, \dots, Y_N} f(Y_1, \dots, Y_N)
           - E_{Y_i, Y_{i+1}, \dots, Y_N} f(Y_1, \dots, Y_N)}              \\
     & = \abs{E_{Y_{i+1}, \dots, Y_N}
           \left[ f(Y_1, \dots, Y_N) - E_{Y_i} f(Y_1, \dots, Y_N) \right]} \\
     & \leq E_{Y_{i+1}, \dots, Y_N}
    \abs{ f(Y_1, \dots, Y_N) - E_{Y_i} f(Y_1, \dots, Y_N) }                \\
     & \leq 1 \text{ a.e.}
  \end{align*}
  Therefore $d_i \leq 1$ and the desired statement follows from the
  previous theorem.

\end{proof}

\begin{eg}
  Consider a discrete cube $\{-1, 1\}^n$ with the Hamming metric, i.e.\
  the distance is the number of coordinates that does not agree.
  Suppose $f : \left\{ -1, 1 \right\}^n \to \R$ is Lipschitz
  with constant $1$. Then, by previous theorems, we have
  \[
    P \left\{ \abs{f(V) - E f(V)} > t \right\}
    \leq 2 \exp \left( - \frac{t^2}{4 n} \right)
  \]
  for all $t > 0$, where $V$ is the uniform random distribution
  on $\{ -1, 1 \}^n$.

  Moreover, for $A \subset \{-1, 1\}^n$ with $\abs{A} = 2^{n - 1}$.
  We have the following \emph{Hamming concentration inequality}:
  \[
    P \left\{ \dist(X, A) \geq t \right\}
    \leq 2 \exp \left( - \frac{C t^2}{n} \right)
  \]
  for some universal constant $C$.
  From this we can obtain $E [\dist (A, V)] = O(\sqrt{n})$.
\end{eg}

Next we present a combinatorial application of martingale method.
\begin{eg}
  Let $G \sim G(n, \frac{1}{2})$ be a uniform distribution on labeled
  $n$-vertex graph. That is, the graph is constructed with
  $\binom{n}{2}$ independent edges $(b_{ij})_{i > j}$ following
  $\Bernoulli(\frac{1}{2})$ distribution.
  We want to consider the chromatic number $\chi (G)$ of the graph.

  Define a filtration by $\fcal_m = \sigma(b_{ij}: j < i \leq m)$,
  $\fcal_0 = \{ \emptyset, \Omega \}$. Now the filtration
  $\fcal_0 \subset \fcal_1 \subset \cdots \subset \fcal_n$ has
  $n$ elements. Now $\seqm{E[\chi(G) | \fcal_m]}^n$ is a martingale,
  so by previous theorems,
  \[
    P \left\{ \abs{\chi (G) - E \chi(G)} \geq t \right\}
    \leq 2 \exp \left(  - \frac{t^2}{4 \sumi^n d_i^2} \right),
  \]
  where $d_i = \norm{E [\chi(G) | \fcal_i] - E [\chi(G) | \fcal_{i - 1}]}_\infty$.
  By intergrating out the extra variables (similar to previous
  example), we have
  \begin{align*}
    E[\chi(G) | \fcal_{l - 1}] = E[E_{b_{lj} : l > j} [\chi(G)] | \fcal_l].
  \end{align*}
  Hence,
  \[
    d_l = \norm{E [\chi(G) -  E_{b_{lj} : l > j} [\chi(G)] | \fcal_l]}_\infty
  \]

  Now for any graph $\Gamma$ on $[n]$, let $T(\Gamma)$ be the set of all
  graphs $\tilde{\Gamma}$ on $[n]$ such that $\{i, j\}$, $i > j$ and
  $i \neq l$ is an edge of $\tilde{\Gamma}$ iff $\{i, j\}$ is
  an edge of $\Gamma$. That is, the graphs can only disagree on edges
  incident from $l$. Then we have
  \[
    E_{b_{lj} : l > j} [\chi(G)] = \frac{1}{\abs{T(G)}} \sum_{\tilde{\Gamma}
      \in T(G)} \chi (\tilde{\Gamma}).
  \]
  This implies that $d_l \leq 1$ a.e., since disagreeing on
  one vertex change the chromatic by at most one. Thus $\var \chi(G) = O(n)$
  and $\chi(G) = E [\chi(G)] (1 \pm o(1))$. Note $E [\chi(G)] \sim
    \frac{n}{\log n}$.

\end{eg}

\subsection{Predicatble sequences and stopping time}

\begin{defi}[predicatble sequence]
  Let $\{ \fcal_n \}_{n = 0}$ be a finite or infinite filtration
  and $\{ H_n \}_n$ be a sequence of random variables.
  Say $\{ H_n \}$ is predicatble if for any $n$,
  $H_n$ is $\fcal_{n - 1}$ measurable.
\end{defi}

Now let $X_n$ be an adapted sequence, and define
a sequence $\{ H \cdot X \}_n$ via
\[
  (H \cdot X)_n = \sumi^n H_i (X_i - X_{i - 1})
\]
for $n \geq 0$.

\begin{thm}
  Suppose $\{ H_n \}$ is non-negative predicatble sequence
  and $\{ X_n \}$ is a submartingale (resp.\ supermartingale). Then
  $\{H \cdot X\}_n$ is a submartingale (resp.\ supermartingale).
\end{thm}

\begin{proof}
  Note that
  \begin{align*}
    E \left[ \sumi^n H_i (X_i - X_{i - 1}) \middle| \fcal_{n - 1} \right]
     & = \sumi^{n - 1} H_i (X_i - X_{i - 1})
    + E[H_n (X_n - X_{n - 1}) | \fcal_{n - 1}] \\
     & = \sumi^{n - 1} H_i (X_i - X_{i - 1})
    + H_n E[X_n - X_{n - 1} | \fcal_{n - 1}].
  \end{align*}
  The statement then follows immediately.
\end{proof}

The main applications of predicatble sequences are related
to stopping times.

\begin{defi}[stopping time]
  A random variable $T$ is a stopping time w.r.t.\ filtration
  $\seqinfn{\fcal_n}$ if for any $k \in \N$,
  the event $\{ T \leq k \}$ is $\fcal_k$-measurable.
\end{defi}

\begin{eg}
  Let $\zseqinfn{X_n}$ be a martingale w.r.t\
  filtration $\zseqinfn{\fcal_n}$ with $\fcal_0$ being the trivial
  $\sigma$-field. Let $T$ be
  a stopping time w.r.t.\ $\zseqinfn{\fcal_n}$
  such that there exists $M < \infty$ satisfying $P \{ T \leq M \} = 1$.
  Then $E X_T = X_0$.

  We can write
  \[
    X_T = X_0 + \sumi^M \ind_{\{T \geq i\}} (X_i - X_{i - 1}).
  \]
  Now for $i \geq 0$,
  define $H_i = \ind_{\{T \geq i\}} = 1 - \ind_{\{T \leq i - 1\}}$,
  so $H_i$ is a predicatble sequence. By the previous theorem,
  the statement follows.
\end{eg}

Next we consider the important optional stopping time theorem.

\begin{thm}[Optional stopping time theorem]
  Let $\zseqinfn{\fcal_n}$ be a filtration, $T$ be a stopping time,
  and $\zseqinfn{X_n}$ an adapted sequence. Then
  if $\zseqinfn{X_n}$ is a submartingale (resp.\ supermartingale), then
  $\zseqinfn{X_{n \wedge T}}$ is also a submartingale (resp.\ supermartingale)
  w.r.t\ the same filtration $\zseqinfn{\fcal_n}$. Moreover, for any fixed
  $M \geq 0$, the sequence $\zseqinfn{X_{(n \vee T) \wedge M}}$
  is a submartingale (resp.\ supermartingale) w.r.t\ the filtration
  $\zseqinfm{\sigma (X_{(j \vee T) \wedge M}, 0 \leq j \leq m)}$.
\end{thm}

\begin{proof}
  For the first part, we have
  \[
    X_{T \wedge n} - X_0 = \sumk^n \ind_{\{T \geq k\}} (X_k - X_{k - 1}).
  \]
  Since $\ind_{T \leq k}$ is a predicatble sequences, $X_{T \wedge n}$
  is a submartingale.

  For the second part, write $\gcal_{m} = \sigma (X_{(j \vee T) \wedge M},
    0 \leq j \leq m)$. Fix any $A \in \gcal_{n-1}$ with $P(A) > 0$ and
  any $k \in \N \cup \{\infty\}$ such that $P(A \cap \{T = k\}) > 0$.
  The event $A$ is representable as
  $\{ (X_{M \wedge (m \vee T)})_{m=0}^{n - 1} \in B \}$ for some
  $B \in \bcal(\R^n)$. Next consider
  \begin{align*}
    E \left[ \left( X_{M \wedge (T \vee n)} - X_{M \wedge (T \vee (n - 1))} \right)
        \ind_{ A \cap \{T = k\} } \right].
  \end{align*}
  There are several cases.

  If $k \geq n$, then $T \vee n = T \vee (n - 1) = T = k$
  on $\{T = k\}$, and the conditional expectation is zero.

  If $k \leq n - 1$, then $T \vee n = n$ and $T \vee (n - 1) = n - 1$.
  The conditional expectation becomes
  \[
    E \left[ \left( X_{M \wedge n} - X_{M \wedge (n - 1)} \right)
      \ind_{ A \cap \{T = k\} } \right].
  \]
  If $M \leq n - 1$, then this becomes zero too. Therefore the only
  nontrivial case is when $M \geq n$. For this case, the conditional
  expectation becomes $E [X_n - X_{n - 1} \mid A \cap \{ T = k \}]$.
  Note that on the event $\{ T = k \}$ with $k \leq n - 1$,
  for all $1 \leq j \leq n - 1$ we have
  \[
    (j \vee T) \wedge M \leq (j \vee (n - 1)) \wedge M \leq n - 1.
  \]
  Therefore, $A$ is
  $\fcal_{n - 1}$-measurable. Also, $T$ is a stopping time, so
  $\{ T = k \}$ is also $\fcal_{n - 1}$-measurable. Now
  submartingale property implies
  \begin{align*}
    E [\left( X_n - X_{n - 1} \right) \ind_{A \cap \{ T = k \}}]
    = E [\ind_{A \cap \{T = k\}} E[X_n - X_{n - 1} | \fcal_{n - 1}]]
    \geq 0.
  \end{align*}
  Therefore,
  \[
    E \left[ \left( X_{M \wedge (T \vee n)} - X_{M \wedge (T \vee (n - 1))} \right)
      \ind_{A \cap \{T = k\}} \right] \geq 0.
  \]
  Now since $A = \bigsqcup_{k \in \N \cup \{\infty\}} A \cap \{T = k\}$,
  we have
  \begin{align*}
    E \left[ \left( X_{M \wedge (T \vee n)} - X_{M \wedge (T \vee (n - 1))} \right)
        \ind_A \right]
     & = \sum_{k \in \N \cup \{\infty\}}
    E \left[ \ind_{A \cap \{T = k\}} (X_{M \wedge (T \vee n)} - X_{M \wedge (T \vee (n - 1))})
        \right] \\
     & \geq 0.
  \end{align*}
  for any $A \in \gcal_{n - 1}$. This then implies that
  \[
    E \left[ X_{M \wedge (T \vee n)} - X_{M \wedge (T \vee (n - 1))}
      | \gcal_{n - 1} \right] \geq 0,
  \]
  as desired.
\end{proof}

\begin{thm}[maximal inequalities]
  Let $\zseqinfn{X_n}$ be a submartingale. Then,
  \[
    a P \left\{ \max_{0 \leq i \leq n} X_i > a \right\}
    \leq E \left[ X_n \ind_{\{\max_{0 \leq i \leq n} X_i > a \}} \right]
    \leq E \abs{X_n}
  \]
  for any $a > 0$.
\end{thm}

\begin{proof}
  Let stopping time $T = \inf \left\{ m \geq 0 : X_m > a \right\}$.
  Then,
  \begin{align*}
    a P \left\{ \max_{0 \leq i \leq n} X_i > a \right\}
     & \leq E \left[ \ind_{\{ \max_{0 \leq i \leq n}
                  X_i > a \}} X_{T \wedge n} \right]                \\
     & = E X_{T \wedge n} -
    E [\ind_{\{\max_{0 \leq i \leq n} X_i \leq a\}} X_{T \wedge n}] \\
     & = E X_{T \wedge n} -
    E [\ind_{\{\max_{0 \leq i \leq n} X_i \leq a\}} X_{n}]          \\
  \end{align*}
  By optional stopping time theorem, $\seqinfm{X_{n \wedge (T \vee m)}}$
  is a submartingale. Then the expectation of
  the $n$-th element is at least the expectation of
  the $0$-th element. This implies $E X_{T \wedge n} \leq E X_n$.
  Alternatively, we can also see this from
  \[
    X_n - X_{n \wedge T} = \sumi^n \ind_{\{T < i\}} (X_i - X_{i - 1}).
  \]
  Rearranging the inequality gives the desired result.

\end{proof}

\begin{cor}[Kolmogorov's maximal inequality]
  Let $\seqinfn{X_n}$ be mutually independent, $E X_n = 0$,
  and $\var X_n < \infty$. Then,
  \[
    P \left\{ \max_{1 \leq n \leq m} \abs{\sumi^n X_i} > a \right\}
    \leq \frac{1}{a^2} \sumn^m \var X_n
  \]
  for any $a > 0$.
\end{cor}

\begin{proof}
  Let $Y_n = \left( \sumi^n X_i \right)^2$ for $n \geq 1$.
  We can verify that $Y_n$ is a submartingale. It then
  follows from the previous theorem that
  \[
    P \left\{ \max_{1 \leq n \leq m} Y_n > a^2 \right\}
    \leq \frac{1}{a^2} E \abs{Y_m} = \frac{1}{a^2} \sumn^m \var X_n,
  \]
  as desired.
\end{proof}

Without martingale theory, we can also provide a more elementary
proof for Kolmogorov's inequality for symmetric variables.

\begin{thm}
  Let $\seqinfn{X_n}$ be mutually independent and symmetric random
  variables. Suppose also $E X_n = 0$,
  and $\var X_n < \infty$. Then,
  \[
    P \left\{ \max_{1 \leq n \leq m} \sumi^n X_i > a \right\}
    \leq \frac{1}{a^2} \sumn^m \var X_n
  \]
  for any $a > 0$.
\end{thm}

\begin{proof}
  Write $S_n = \sumi^n X_i$. We have
  \begin{align*}
    P \left\{ \max_{1 \leq i \leq m} S_i > a \right\}
     & \leq P \left\{ S_m > a \right\}
    + P \left\{ \exists i < m : S_i > a \text{ and } S_m < S_i \right\} \\
     & = P \left\{ S_m > a \right\}
    + \sumi^{m - 1} P \left\{ S_i > a, \; S_r \leq a
    \text{ for all $r < i$, and } S_m - S_i < 0 \right\}                \\
     & = P \left\{ S_m > a \right\}
    + \sumi^{m - 1} P \left\{ S_i > a, \; S_r \leq a
    \text{ for all $r < i$, and } S_m - S_i > 0 \right\},
  \end{align*}
  where in the last step we used the fact that $X_i$'s are symmetric.
  Then, combining $S_i > a$ and $S_m - S_i > 0$, we can obtain
  \[
    P \left\{ \max S_i > a \right\}
    \leq 2 P \left\{ S_m > a \right\}
    = P \left\{ \abs{S_m} > a \right\}
    \leq \frac{1}{a^2} \sumi^m \var X_i.
  \]

\end{proof}

\subsection{Martingale convergence}

\begin{thm}[upcrossing lemma]
  Let $\zseqinfn{X_n}$ be a submartingale and fix $a < b$.
  Next define stopping times $\zseqinfk{T_k}$ taking values
  in $\N \cup \{\infty, 0, -1\}$ via $T_0 = -1$, and for
  $k \geq 1$,
  \[
    T_{2k - 1} = \inf \{l > T_{2k-2} : X_l \leq a\}, \quad
    T_{2k} = \inf \{l > T_{2k-1}: X_l \geq b\}.
  \]
  Also, for all $m \geq 0$, let $U_m = \max \{k \geq 0:  T_{2k} \leq m \}$.
  Then, for any $m \geq 0$,
  \[
    (b - a) E [U_m] \leq E [(X_m - a)^+] - E [(X_0 - a)^+].
  \]
\end{thm}

\begin{remark}
  If $X_n$ is a submartingale, then $(X_n - a)^+$ is also a
  submartingale. Indeed, since $\phi(x) = \max(x, 0)$ is convex and
  monotone,
  \[
    E [\phi(X_{m + 1}) | \fcal_m] \geq \phi (E [X_{m + 1} | \fcal_m])
    \geq \phi(X_m).
  \]
\end{remark}

\begin{proof}
  For $n \geq 1$, define
  \[
    H_n = \ind_{\{T_{2k-1} < n \leq T_{2k} \text{ for some $k \geq 1$}\}}.
  \]
  Claim that $H_n$ is predicatble. Indeed, note that
  $H_n = 1$ iff $\max \{l : T_l < n\}$ is odd, and this is
  $\fcal_{n - 1}$ measurable. Write $Z_n = (X_n - a)^+$.
  Claim that
  \begin{align*}
    (b - a) U_m
    \leq \sum_{i=1}^{\max \{T_{2l} : T_{2l} \leq m\}}
    H_i (Z_i - Z_{i - 1})
    = (H \cdot Z)_{\max \{T_{2l} : T_{2l} \leq m\}}.
  \end{align*}
  Indeed, each upcrossing contributes at least $b - a$ to the summation.
  By the condition of $H_n = 1$ above, we can verify that
  \[
    \sum_{i = \max \{T_{2l} : T_{2l} \leq m\} + 1}^m H_i (Z_i - Z_{i - 1}) \geq 0.
  \]
  Hence, $(b - a) U_m \geq (H \cdot Z)_m$. Now note that $Z$ is a
  submartingale and $1 - H$ is predicatble, so
  $((1 - H) \cdot Z)$ is also a submartingale. Therefore,
  \begin{align*}
    (b - a) E[U_m]
     & \leq E (H \cdot Z)_m                         \\
     & \leq E (H \cdot Z)_m + E ((1 - H) \cdot Z)_m \\
     & = E (1 \cdot Z)_m                            \\
     & = E Z_m - E Z_0.
  \end{align*}
\end{proof}

Now an application of the upcrossing lemma.

\begin{eg}
  Suppose $E [(X_n - a)^+]$ is uniformly bounded, i.e.\
  \[
    \sup_n E [(X_n - a)^+] < \infty.
  \]
  This implies that $E U_m$ is also uniformly bounded.

  Consider the event $A$ that there is infinitely many oscillations.
  Then everywhere on $A$, we have $\lim_{m \to \infty} U_m = \infty$.
  Hence if $P(A) > 0$, then $E U_m \to \infty$ as $m \to \infty$, a
  contradiction.
\end{eg}

We can also use this to consider a symmetric random walk on
the real numbers.

\begin{eg}
  Fix $0 < \gamma < 1$.
  Let $\seqinfn{Y_n}$ be mutually independent and symmetric
  random variable with
  \[
    P \left\{ Y_n = n^{-\gamma} \right\}
    = P \left\{ Y_n = - n^{-\gamma} \right\}
    = \frac{1}{2}.
  \]
  Let $X_n = \sumj^n Y_j$. Then $X_n$ is a martingale and
  we hope to investigate the trajectory of $X_n$.

  We have
  \[
    \var X_n = \sumk^n \frac{1}{k^{2 \gamma}}.
  \]
  Hence, if $\gamma > \frac{1}{2}$, the variance of $X_n$ are
  uniformly bounded. Therefore, for any $a \in \R$,
  $E [(X_n - a)^+]$ is also uniformly bounded, and thus
  by our previous example, for any $a < b$,
  there are finitely many oscillations with probability $1$.
  From the previous example and some extra argument,
  we can derive that the random walk converges to a limit that is nonzero.
  Indeed, the variance is nonzero, so the limit is not a.s.\ equal to zero.
  From Kolmogorov $0$-$1$ law, the limit is nonzero a.s.
\end{eg}

\begin{thm}[submartingale convergence theorem]
  Let $\seqinfn{X_n}$ is a submartingale with
  \[
    \sup_n E[X_n^+] < \infty.
  \]
  Then there exists random variable $X$ such that
  $E \abs{X} < \infty$ and $X_n \to X$ a.e.
\end{thm}

\begin{proof}
  Let $S = \limsup_n X_n$ and $L = \liminf_n X_n$.
  Then we have $L^+ \leq \liminf_n X_n^+$ and it follows from
  Fatou's lemma that
  \[
    E [L^+] \leq E[\liminf_n X_n^+] \leq \liminf_n E[X_n^+] < \infty
  \]
  Also,
  \[
    \sup_n E[(-X_n)^+] = \sup_n \left( E [X_n^+] - E X_n \right)
    \leq \sup_n E[X_n^+] - E X_1
    < \infty.
  \]
  Now $-S = \liminf_n (- X_n)$ so it follows similarly from
  Fatou's lemma that $E [(- S)^+] < \infty$. Therefore,
  $S \neq \infty$ and $L \neq - \infty$ a.e.

  Now suppose $P \{ S - L > \epsilon \} > 0$.
  Then there exists $a < b$ such that $P \{ L < a < b <
    S \} > 0$. This implies that there is a positive probability
  that infinitely oscillations happens between $a, b$ and thus
  $E U_m \to \infty$ as $m \to \infty$, a
  contradiction with $\sup_n E[X_n^+] < \infty$. Moreover,
  $E \abs{X} < \infty$ since $E [X^+] = E [L^+] < \infty$
  and $E [X^-] = E [(-S)^+] < \infty$.

\end{proof}

\begin{cor}
  Let $\seqinfn{X_n}$ be a non-negative supermartingale.
  Then there exists random variable $X$ such that
  $E X \leq E X_n$ for any $n \geq 0$ and $X_n \to X$ a.e.
\end{cor}

\begin{proof}
  Note that $-X_n$ is a submartingale and $-X_n \leq 0$.
  Hence $\sup_n E [(-X_n)^+] = 0$, so existence of $X$ is guaranteed
  by the previous theorem. Moreover,
  \[
    E[X] \leq \liminf_{n \to \infty} E[X_n] = \inf_n E[X_n]
  \]
  since $X_n$ is a supermartingale, completing the proof.

\end{proof}

\begin{eg}
  Let $\xi_1, \xi_2, \dots$ i.i.d.\ non-negative random variables
  with $E \xi_1 \leq 1$. Then $X_n = \prod_{i=1}^n \xi_i$
  is a supermartingale, and $X_n \to 0$ a.e.\ unless $\xi_1 = 1$ a.e.

  To prove this,
  suppose $\xi_1 \not\equiv 1$ then there exists $\epsilon > 0$ such that
  $P \left\{ \abs{\xi_1 - 1} > \epsilon \right\} > \epsilon$.
  Then, the event $A = \{ \abs{\xi_i - 1} \geq \epsilon \text{ i.o.} \}$
  has probability one.
  On event $A$,
  \[
    \abs{X_{m + 1} - X_m} =
    \abs{\prod_{i=1}^{m + 1} \xi_i - \prod_{i=1}^m \xi_i} \geq
    \epsilon \prod_{i=1}^m \xi_i
  \]
  for all $m$ such that $\abs{\xi_{m+1} - 1} \geq \epsilon$,
  and i.e.\ for infinitely many indices. By convergence theorem,
  there exists random variable $X$ such that $X_n \to X$ a.e.
  Let $B$ be the event $X_n \to X$.

  To show $X = 0$, suppose for contradiction that
  there exist $\delta > 0$ such that $ P \left\{ X > \delta \right\}
    > \delta$, then let event $C = \{ X > \delta \}$.
  Consider the intersection $A \cap B \cap C$.
  Let $\omega \in A \cap B \cap C$. Then for all sufficient
  large $n$, $X_n (\omega) > \delta$. Then for infinitely many $m$'s
  we have $\abs{X_{m + 1} (\omega) - X_{m} (\omega)} \geq \epsilon \delta$.
  Therefore, $\lim_{n \to \infty} X_n (\omega)$ does not exist, a
  contradiction.

\end{eg}

\begin{eg}[discrete-time birth process]
  \def\tx{\tilde{X}}
  Let $\{ \xi_{g, i} \}_{g, i = 1}^\infty$ be i.i.d.\ non-negative
  integer valued random variables. Define $\seqinfn{X_n}$ inductively
  by $X_0 = 1$ and $X_n = \sum_{i=1}^{X_{n-1}} \xi_{n, i}$ for $n \geq 1$.
  We hope to investigate $P \{ x_n > 0 \text{ for all $n \geq 1$} \}$.
  Claim that
  \begin{itemize}
    \item $E \xi_{1, 1} > 1$ implies $P \{ x_n > 0 \text{ for all $n \geq 1$} \} > 0$.
    \item $E \xi_{1, 1} < 1$ implies $P \{ x_n > 0 \text{ for all $n \geq 1$} \} = 0$.
    \item $E \xi_{1, 1} = 1$ and $\xi_{1, 1} \not\equiv 1$
          implies $P \{ x_n > 0 \text{ for all $n \geq 1$} \} = 0$.
  \end{itemize}

  First consider the caes where $E \xi_{1, 1} < 1$. Claim that
  $X_n$ is non-negative supermartingale. Indeed,
  \begin{align*}
    E [X_n | \fcal_{n - 1}]
     & = E \left[ \sum_{i=1}^{X_{n-1}} \xi_{n, i} \middle| \fcal_{n - 1} \right] \\
     & = E \left[ \suminfi \xi_{n, i} \ind_{\{i \leq X_{n - 1}\}}
             \middle| \fcal_{n - 1} \right]                 \\
     & = \suminfi \ind_{\{i \leq X_{n - 1}\}} E \xi_{1, 1}                       \\
     & \leq X_{n - 1}.
  \end{align*}
  Then there exists non-negative $X$ such that $X_n \to X$ a.e.
  Moreover, it is easy to verify that $\lim_{n \to \infty} E X_n = 0$.
  Hence by Fatou's lemma,
  $0 = \lim_{n \to \infty} E X_n \geq E \lim_{n \to \infty} X_n$,
  which implies $E X = 0$ and thus $X = 0$ a.e.

  Next consider the case where $E \xi_{1,1} = 1$ with $\xi_{1,1}$
  non-constant. We still have that $X_n$ is a non-negative
  supermartingale. Therefore, there exists $X \geq 0$ such that
  $X_n \to X$ a.e. It is easy to see that $X$ is integer-valued.
  Suppose there exists $m \geq 1$ such that $P \{ X = m \} > 0$.
  On the event $\{ X = m \}$, $X_n = m$ for all but finitely many $n$'s.
  However, since $\xi_{1,1} \not \equiv 1$, $X_n$ does not
  a.s.\ equals $m$ given $X_{n-1} = m$. Therefore, for
  some $\delta > 0$,
  \begin{align*}
    P \left\{ X_n = m \text{ for all $n_0 \leq n \leq N$} \right\}
     & = \prod_{n = n_0 + 1}^N P \left\{ X_n = m | X_j = m
    \text{ for all $n_0 \leq j \leq n$} \right\}           \\
     & \leq (1 - \delta)^{N - n_0},
  \end{align*}
  a contradiction.

  Finally, consider the last case where $E \xi_{1,1} > 1$.
  We first introduce the concept of \emph{stochastic domination}.
  We say random variable $Z$ \emph{stochastically dominates}
  $W$ if $P \{W \leq t\} \geq P \{Z \leq t\}$ for all $t \in \R$.
  In this case, there exists $\tilde{W} \sim W$ and $\tilde{Z} \sim Z$
  on the same probability space such that $\tilde{W} \leq \tilde{Z}$ a.e.

  Let $M > 0$ and $\epsilon > 0$ be such that $\tilde{\xi}_{1,1}
    = \min (\xi_{1,1}, M)$ satisfy $E \tilde{\xi}_{1,1} \geq 1 + \epsilon$.
  Also define an auxilary process via $\tilde{X}_0 = 1$
  and $\tilde{X}_n = \sumi^{\tilde{X}_{n - 1}} \tilde{\xi}_{n, i}$.
  We can verify that $X_n$ stochastically dominates
  $\tilde{X}_n$ and $\tilde{X}_n \leq X_n$ a.e.\ for any $n$.

  Claim that $P \{ \tilde{X}_n > 0 \text{ for all } n \geq 1 \} > 0$.
  Write $\fcal_n = \sigma (\tx_1, \dots, \tx_{n})$.
  Since $E \tilde{\xi}_{1,1} \geq 1 + \epsilon$, we have
  $E [\tilde{X}_n | \tx_{n - 1}] \geq (1 + \epsilon) \tilde{X}_{n - 1}$
  a.e. Hence,
  \begin{align*}
    P \left\{ \tilde{X}_n \leq \left( 1 + \frac{\epsilon}{2}
    \tilde{X}_{n - 1} \right) \middle| \fcal_{n - 1} \right\}
     & \leq P \left\{ \tilde{X}_n \leq \frac{1 + \frac{\epsilon}{2}}
                                       {1 + \epsilon} E [\tilde{X}_n | \tx_{n - 1}]
    \middle| \fcal_{n - 1} \right\}                                                 \\
     & \leq P \left\{ \tx_n \leq \left( 1 - \frac{\epsilon}{4} \right)
    E [\tx_n | \tx_{n - 1}] \middle| \fcal_{n - 1} \right\},
  \end{align*}
  where in the last step we chose $\epsilon > 0$ to be small
  enough such that
  \[
    \frac{1 + \frac{\epsilon}{2}}{1 + \epsilon} \leq 1 - \frac{\epsilon}{4}.
  \]
  Write $A_k = \{ \tx_{n - 1} = k \}$ with $P(A_k) > 0$, then
  for $k > 0$ we have
  \begin{align*}
    P \left\{ \tx_n \leq \left( 1 - \frac{\epsilon}{4} \right)
    E [\tx_n | \tx_{n - 1} = k ] \middle| \tx_{n - 1} = k \right\}
     & \leq \frac{\var [\tx_n | \tx_{n - 1} = k]}{(\frac{\epsilon}{4})^2
              E[\tx_n | \tx_{n - 1} = k]^2}.
  \end{align*}
  For $k = 0$ simply bound the probability with $1$. Now,
  \[
    \var [\tx_n | \tx_{n - 1} = k]
    = k \var \tilde{\xi}_{1,1} \leq k M^2
  \]
  and
  \[
    E [\tx_n | \tx_{n - 1} = k] = k E \tilde{\xi}_{1,1}
    \geq k.
  \]
  Therefore,
  \begin{align*}
    P \left\{ \tx_n \leq \left( 1 - \frac{\epsilon}{4} \right)
    E [\tx_n | \tx_{n - 1} = k ] \middle| \tx_{n - 1} = k \right\}
     & \leq \frac{16 M^2}{\epsilon^2 k}
  \end{align*}
  and thus
  \[
    P \left\{ \tilde{X}_n \leq \left( 1 + \frac{\epsilon}{2}
    \tilde{X}_{n - 1} \right) \middle| \fcal_{n - 1} \right\}
    \leq \frac{16 M^2}{\epsilon^2 \tx_{n - 1}}
  \]
  Define $n_0 = \ceil{64 M^2 \epsilon^{-3}}$,
  and event $B_n = \left\{ \tx_{n} \geq \left( 1 + \frac{\epsilon}{2} \right)^{n - n_0}
    64 M^2 \epsilon^{-3} \right\}$ for $n \geq n_0$.
  Since $E \tilde{\xi}_{1,1} > 1$, $P(B_{n_0}) > 0$. Also
  It follows that
  \[
    P (B_{n_0 + 1} | B_{n_0}) \geq 1 - \frac{16 M^2}{\epsilon^2
      \cdot 64 M^2 \epsilon^{-3}} = 1 - \frac{\epsilon}{4}.
  \]
  Similarly, we have $P(B_{n_0 + 2} | B_{n_0 + 1} \cap B_{n_0})
    \geq 1 - \frac{1}{4} (1 + \frac{\epsilon}{2})^{-1} \epsilon$,
  etc. Therefore,
  for the event $C = \left\{ \tx_n \geq \left( 1 + \frac{\epsilon}{2}
    \right)^{n - n_0} \tx_{n_0}
    \text{ for all } n > n_0 \text{ and }
    \tx_{n_0} \geq 64 M^2 \epsilon^{-3} \right\}$, we have
  \begin{align*}
    P (C) \geq P \left\{ \tx_{n_0} \geq 64 M^2 \epsilon^{-3} \right\}
    \prod_{n = n_0 + 1}^\infty \left( 1 -
                               \frac{1}{4} \left( 1 + \frac{\epsilon}{2} \right)^{n_0 - n}
    \epsilon \right) > 0.
  \end{align*}

\end{eg}


\subsection{Unifrom integrability, convergence in $L^1$}

\begin{thm}[uniform integrability]
  Let $\seqinfn{X_n}$ are random variables with $E \abs{X_n} < \infty$.
  Suppose $X$ is another random variable such that $X_n$ converges
  to $X$ in probability. Then the following are equivalent:
  \begin{enumerate}
    \item $X_n \to X$ in $L^1$.
    \item $\seqinfn{X_n}$ is uniformly integrable, i.e.\
          \[
            \lim_{M \to \infty} \sup_n E \left[ \abs{X_n} \ind_{\{\abs{X_n} > M\}}
              \right] = 0.
          \]
  \end{enumerate}
\end{thm}

\begin{proof}
  (1) $\implies$ (2). Suppose $X_n \to X$ in $L^1$, then $E \abs{X} < \infty$
  and $s = \sup_n E \abs{X_n} < \infty$. Fix $\epsilon > 0$, then
  there exists $\delta > 0$ such that $A \in \Sigma$
  with $P (A) \leq \delta$ implies $E [\abs{X} \ind_A] < \epsilon$.
  Let $n_0$ be such that $n \geq n_0$ implies
  $E \abs{X - X_n} \leq \epsilon$. There also exists $M_0 < \infty$
  such that
  \[
    \max_{n < n_0} E \left[ \abs{X_n} \ind_{\{\abs{X_n} \geq M_0 \}} \right]
    < \epsilon.
  \]
  Now let $M = \max (M_0, \delta^{-1} s)$. Then for $n \geq n_0$,
  we have
  \begin{align*}
    E \left[ \abs{X_n} \ind_{\{\abs{X_n} \geq M\}} \right]
    \leq E \abs{X - X_n} + E \left[ \abs{X} \ind_{\{\abs{X_n} \geq M\}} \right]
    \leq 2 \epsilon,
  \end{align*}
  where for the second term we used Markov inequality
  to obtain $P \{\abs{X_n} \geq M\} \leq M^{-1} s \leq \delta$.

  (2) $\implies$ (1). Suppose $X_n$ does not converge to $X$
  in $L^1$. Then there is a subsequence $\seqinfk{X_{n(k)}}$ and
  $\epsilon > 0$ such that $E \abs{X_{n(k)} - X} \geq \epsilon$ for
  all $k$. To simplify notation without loss of generality,
  consider $X_n$ with $E \abs{X_n - X} \geq \epsilon$ for all $n$.

  In the case where $E \abs{X} = \infty$, with $X_n \convp X$
  claim that $\sup_n E \abs{X_n} = \infty$. Indeed,
  since $X_n \to X$ in probability, there is a subsequence $X_{n(k)}$
  such that $X_{n(k)} \to X$ a.s. If $\sup_n E \abs{X_n}$
  is finite, Fatou's lemma then gives
  \[
    E \abs{X} \leq \liminf E \abs{X_n} < \infty,
  \]
  a contradiction. Therefore, $\sup_n E \abs{X_n} = \infty$
  and $X_n$ is not uniformly integrable.

  Now suppose $E \abs{X} < \infty$. Fix $\epsilon > 0$.
  There is $\delta > 0$ such that $P(A) \leq \delta$ implies
  $E \left[ \abs{X_n} \ind_A \right] \leq \frac{\epsilon}{4}$.
  Since $X_n - X \convp 0$, we have
  $P \left\{ \abs{X_n - X} \geq 1 \right\} \to 0$
  and thus
  \[
    E \left[ \abs{X_n - X} \ind_{\{\abs{X_n - X} < 1\}}\right]
    \leq P \left\{ \abs{X_n - X} \geq 1 \right\} \to 0
  \]
  as $n \to \infty$.
  Now for $n$ large enough,
  \begin{align*}
    E \left[ \abs{X_n} \ind_{\{\abs{X - X_n} \geq 1\}} \right] + \frac{\epsilon}{4}
     & \geq E \left[ \abs{X_n} \ind_{\{\abs{X - X_n} \geq 1\}} \right]
    + E \left[ \abs{X} \ind_{\{\abs{X - X_n} \geq 1\}} \right]             \\
     & \geq E \left[ \abs{X - X_n} \ind_{\{\abs{X_n - X} \geq 1\}} \right] \\
     & \geq E \abs{X - X_n} - \frac{\epsilon}{4}                           \\
     & \geq \frac{3 \epsilon}{4}.
  \end{align*}
  Let $P \left\{ \abs{X_n - X} \geq 1 \right\} = \delta_n$, then
  $\delta_n \to 0$. We then have
  \begin{align*}
    \liminf_{n \to \infty} E \left[ \abs{X_n}
                               \ind_{\{\abs{X_n} > \delta_n^{-1/2}\}} \right]
     & \geq \liminf_{n \to \infty} E \left[ \abs{X_n} \ind_{\{\abs{X_n} > \delta_n^{-1/2}\}}
                                       \ind_{\{\abs{X_n - X} \geq 1\}} \right].
  \end{align*}
  However,
  \begin{align*}
     & E \left[ \abs{X_n} \ind_{\{\abs{X_n} > \delta_n^{-1/2}\}}
           \ind_{\{\abs{X_n - X} \geq 1\}} \right]     \\
     & = E \left[ \abs{X_n} \ind_{\{\abs{X - X_n} \geq 1\}} \right]
    - E \left[ \abs{X_n} \ind_{\{\abs{X_n} \leq \delta_n^{-1/2}\}}
          \ind_{\{\abs{X_n - X} \geq 1\}} \right]   \\
     & \geq \frac{\epsilon}{2} - \delta_n^{1/2}.
  \end{align*}
  Therefore, $\liminf_n E \left[ \abs{X_n} \ind_{\{\abs{X_n}
        > \delta_n^{-1/2}\}} \right] \geq \frac{\epsilon}{2}$
  and thus $\seqinfn{X_n}$ is not uniformly integrable.

\end{proof}

\begin{thm}
  Let $X$ be a random variable with $E \abs{X} < \infty$ and
  $\seqinfn{\fcal_n}$ be a filtration. Then, the sequence
  $\seqinfn{E [X | \fcal_n]}$ is uniformly integrable.
\end{thm}

\begin{proof}
  Fix $\epsilon > 0$. Let $\delta > 0$ be such that for any event $A$
  with $P(A) \leq \delta$, we have $E [\abs{X} \ind_A] < \epsilon$.
  Also, let $M = \delta^{-1} E \abs{X}$. Claim that for
  any $n$, we have
  \[
    E \left[ \abs{E [X | \fcal_n]} \ind_{\{\abs{E[X | \fcal_n]} \geq M\}}\right]
    \leq \epsilon.
  \]
  Write event $A_M = \{\abs{E[X | \fcal_n]} \geq M\}$. First of all,
  from Jensen's inequality, we have
  \[
    \abs{E [X | \fcal_n]} \ind_{A_M}
    \leq E[\abs{X} | \fcal_n] \ind_{A_M}
  \]
  Note that $\ind_{A_M}$ is $\fcal_n$-measurable, so
  \[
    E[ E[\abs{X} | \fcal_n] \ind_{A_M} ]
    = E[ E[\abs{X} \ind_{A_M} | \fcal_n] ]
    = E [\abs{X} \ind_{A_M}].
  \]
  With Markov's inequality, we know $P(A_M) \leq \delta$.
  By the way $\delta$ is chosen, the proof is complete.

\end{proof}

\begin{thm}
  Let $X$ be a random variable with $E \abs{X} < \infty$ and
  $\seqinfn{\fcal_n}$ be a filtration.
  Also write $\fcal_\infty = \sigma (\fcal_n : n \geq 1)$.
  Then $E[X | \fcal_n]$ converges to $X$ a.e.\ and in $L^1$.
\end{thm}

\begin{proof}
  Towards applying submartingale convergence theorem, we have
  \begin{align*}
    E \left[ E[X | \fcal_n]^+ \right]
    \leq E [ \abs{ E[ X | \fcal_n] } ] \leq E \abs{X}
  \end{align*}
  from Jensen's inequality. By submartingale convergence theorem.
  there exists random variable $Z$ with $E \abs{X} < \infty$
  such that $E [X | \fcal_n] \to Z$ a.e., which implies
  $E [X | \fcal_n] \to Z$ in probability. Together with previous
  theorems, we also have $E [X | \fcal_n] \to Z$ in $L^1$.

  Now the goal is to to show that $Z = E [X | \fcal_\infty]$
  a.e. Let event $A \in \fcal_m$, then for any $n \geq m$,
  $\ind_A E[X | \fcal_n] = E [X \ind_{A} | \fcal_n]$.
  This implies $E [\ind_A E [X | \fcal_n]] = E [\ind_A X]$.
  Now, $E [X | \fcal_n] \to Z$ in $L^1$ implies
  $E [\ind_A E [X | \fcal_n]] \to E [\ind_A Z]$. Therefore,
  $E [\ind_A X] = E [\ind_A Z]$ for any $A \in
    \cupinfm \fcal_m$. Note that $\cupinfm \fcal_m$ is a $\pi$-system
  that generates $\fcal_\infty$. We can also check that
  all sets $A \in \fcal_\infty$ that satisfy
  $E[\ind_A X] = E [\ind_A Z]$ is a $\lambda$-system.
  Applying $\pi$-$\lambda$ theorem and recall the
  definition of conditional expectation compeltes the proof.

\end{proof}

\begin{cor}[Levy $0$-$1$ law]
  Let $\seqinfn{\fcal_n}$ be a filtration and let $\fcal_\infty
    = \sigma (\fcal_n : n \geq 1)$.
  Then, $E [\ind_A | \fcal_n] \to \ind_A$ a.e. for any $A \in \fcal_\infty$.
\end{cor}

\begin{remark}
  This implies Kolmogorov $0$-$1$ law. Indeed, suppose
  $\seqinfn{X_n}$ are mutually independent and let $\tcal =
    \capinfm \sigma \left( X_n : n \geq m \right)$
  be the tail $\sigma$-field. Also let $\fcal_n =
    \sigma (X_1, \dots, X_n)$. Note that
  $\tcal \subset \fcal_\infty$.
  Now take $A \in \tcal$, then
  $A$ is independent of any $\fcal_n$. Therefore
  $E [\ind_A | \fcal_n] = P(A)$. By Levy $0$-$1$ law,
  we can conlude that $P(A) \in \{ 0, 1 \}$.
\end{remark}

\begin{eg}
  Let $\seqinfn{Y_n}$ be mean zero, mutually independent
  random variables. Immediately $\seqinfn{Y_n}$ is a martingale.
  Suppose also $\suminfn E Y_n^2 < \infty$.
  Define $S_n = \sumk^n Y_k$, then
  \[
    E [S_n^+] \leq E \abs{S_n} \leq ( E \abs{S_n}^2 )^{1/2}
  \]
  by Jensen's inequality. Now submartingale convergence theorem
  gives a random variable $X$ with
  $E \abs{X} < \infty$ such that $\suminfk Y_k = X$ a.e.
  As an exercise, we can verify that
  $E [X | Y_1, \dots, Y_n] = \sumk^n Y_k$.
  Now, by previous theorems, we have $\suminfk Y_k \to X$ in $L^1$.
\end{eg}

\subsection{Backward martingales}

\begin{defi}[backward martingale]
  Let $\{\fcal_n\}_{n = -\infty}^0$ be a filtration
  such that $\fcal_{-k} \subset \fcal_{-k+1}$ for all $k \geq 1$.
  Let $\{X_n\}_{n = -\infty}^0$ be adapted to the filtration.
  Suppose $E \abs{X_n} < \infty$ for all $n \leq 0$.
  If $E [X_{n + 1} | \fcal_n] = X_n$ for all $n \leq -1$,
  then say $\{X_n\}_{n = -\infty}^0$ is a \emph{backward martingale}.
\end{defi}

\begin{thm}
  Let $\{X_n\}_{n = -\infty}^0$ be a backward martingale
  under the filtration $\{\fcal_n\}_{n = -\infty}^0$.
  Write $\fcal_{-\infty} = \bigcap_{m = -\infty}^0 \fcal_m$,
  then $X_m \to E[X_0 | \fcal_{-\infty}]$ a.e.\ as $m \to -\infty$.
\end{thm}

\begin{proof}
  For each $m \leq -1$, define sequence $\{ X_{n \wedge 0} \}_{n = m}^\infty$.
  This is a martingale. Fix $a < b$ and let $\{ U_{m, a, b} \}_{m=-1}^{-\infty}$
  be the number of upcrossings. By upcrossing lemma, we have
  \[
    (b - a) E[ U_{m, a, b} ] \leq E [(X_0 - a)^+].
  \]
  We can check that $\{ U_{m, a, b} \}_{m=-1}^{-\infty}$ is monotone
  everywhere on $\Omega$. This allows us to define
  \[
    U_{-\infty, a, b} = \lim_{m \to -\infty} U_{m, a, b}.
  \]
  We have $E[ U_{-\infty, a, b} ]< \infty$. The rest of the argument
  is similar to martingale convergence theorem, and we can show
  that $\liminf_{m \to -\infty} X_m = \limsup_{m \to -\infty} X_m$.
  That is, there exists an a.e.\ limit $\lim_{m \to \infty} X_m =
    X_{-\infty}$ with $E \abs{X_{-\infty}} < \infty$.

  It remains to show that
  $X_{- \infty} = E [X_0 | \fcal_{-\infty}]$.
  With our construction, we have $X_n = E[X_0 | \fcal_n]$.
  This implies that $X_n$'s are uniformly integrable, and thus
  $X_n \to X_{-\infty}$ in $L^1$.
  With $L^1$ convergence, for any event $A$ we have
  $E[\ind_A X_n] \to E[\ind_A X_{-\infty}]$. Note that
  \begin{align*}
    E[\ind_A X_n] = E [\ind_A E[X_0 | \fcal_n]] = E [\ind_A X_0]
  \end{align*}
  for any $A \in \fcal_{-\infty}$. It follows that
  $E[\ind_A X_0] = E[\ind_A X_{-\infty}]$. Since $X_{-\infty}$
  is $\fcal_{-\infty}$-measurable, we can conclude that
  $X_{-\infty} = E[X_0 | \fcal_{-\infty}]$.

\end{proof}


With this, we can revisit law of large numbers.
\begin{thm}[strong law of large numbers]
  Let $X_1, X_2, \dots$ be i.i.d.\ random variables with $E X_1 = 0$.
  Define $S_n = n^{-1} \sumi^n X_i$. We then have
  $S_n \to 0$ a.e.
\end{thm}

\begin{proof}
  Define random variable $Y_{-n} = S_n$ and filtration
  \[
    \fcal_{-n} = \sigma(Y_{-n}, Y_{-n-1}, Y_{-n-2}, \dots) =
    \sigma \left( \sumi^n X_n, X_{n + 1}, X_{n + 2}, \dots \right).
  \]
  That is, we don't have precise information of the first $n$
  elements. Claim that
  $\{ Y_n \}_{n = -\infty}^{-1}$ is a backward martingale, i.e.\
  $E [ Y_{- n + 1} | \fcal_{-n} ] = Y_{-n}$ for all $n \leq -1$.
  This means we need to show that
  \begin{align*}
    E \left[ \frac{1}{n - 1} \sumi^{n - 1} X_i
        \middle| \sumi^n X_i, X_{n + 1}, \dots \right] = \frac{1}{n} \sumi^n X_i.
  \end{align*}
  Let $A \in \sigma (\sumi^n X_i, X_{n + 1}, X_{n + 2}, \dots)$.
  This means $A = \{ (\sumi^n X_i, X_{n + 1}, X_{n + 2}, \dots)
    \in B \}$ for some Borel set $B$. Since $X_i$'s are i.i.d.\
  and $A$ only depends on $X_1, \dots, X_n$ through
  $\sumi^n X_i$, by symmetry we must have
  \[
    P (A \cap \{ X_i \in C \}) = P(A \cap \{ X_j \in C \})
  \]
  for any $1 \leq i,j \leq n$ and any $C \subset \bcal_\R$.
  Therefore, $E [X_i \ind_A]$ is the same for all $1 \leq i \leq n - 1$
  and thus $E [\sumi^{n-1} X_i \ind_A] = (n - 1) E [X_1 \ind_A]$. Similarly,
  $E [\sumi^n X_i \ind_A] = n E[X_1 \ind_A]$. This proves the equation
  and confirms that $Y_n$ is a backward martingale.

  With the previous theorem, we have $Y_{-n} \to E[Y_{-1} | \fcal_{-\infty}]$
  a.e. However, the weak law of large numbers gives $Y_{-n} \to 0$
  in probability, so we must have $Y_{-n} \to 0$ a.e.

\end{proof}

\end{document}